\chapter{Genetics, transcription, translation}
\label{vol06:ch03}

\epigraph{It has not escaped our notice that the specific pairing we
have postulated immediately suggests a possible copying mechanism for
the genetic material.}{James D. Watson and Francis H. C. Crick,
\emph{Nature} 171, 737 (1953), closing paragraph
\cite[p.~738]{paper:watson-crick1953}.}

\chapmeta{Half-life: foundational. Archetypes: information, failure, scaling. Exercise target: 30. Project track: simulation of genetic drift in a finite population (Wright-Fisher model). Hazard class: none (in-silico work). Reader path default: standard, with sections explicitly tagged. Named cases: HeLa cells (1951), He Jiankui CRISPR babies (2018).}

\noindent
A cell that does not replicate dies in one generation. A cell that
replicates without controlled error dies in many generations through
the accumulation of damage. The chapter develops the molecular
machinery by which cells store, copy, transcribe, and translate
genetic information, treats the genetic code as an engineering
specification, quantifies the error rates of replication and repair,
surveys the sequencing technologies that make modern molecular
engineering possible, and closes on the failure modes that arise when
the information machinery breaks (cancer, hereditary disease, ageing).

\section{DNA as digital storage}\pathtag{core}

A cell's genetic information is stored in deoxyribonucleic acid (DNA),
a double-helical polymer of four nucleotides. Each nucleotide consists
of a phosphate group, a deoxyribose sugar, and one of four nitrogenous
bases: adenine (A), thymine (T), guanine (G), or cytosine (C). The
double helix consists of two antiparallel strands held together by
hydrogen bonds between complementary base pairs: A pairs with T
(two hydrogen bonds), G pairs with C (three hydrogen bonds). The
geometry of the helix is set by the sugar-phosphate backbone and the
stacking of the base pairs; the canonical B-form has a diameter of
$2\,\text{nm}$, a rise of $0.34\,\text{nm}$ per base pair, and a
helical pitch of $3.4\,\text{nm}$ (ten base pairs) per turn
\cite{paper:watson-crick1953}. Figure~\ref{fig:dna-geometry}
diagrams the geometry.

\begin{figure}[h]
\centering
% DNA double helix geometry: rise, pitch, diameter, major/minor groove.
\begin{tikzpicture}[>=Stealth, scale=0.95]
  % Two intertwined sinusoids forming the helix backbone
  \def\helixA{0.85}    % amplitude (radius proxy)
  \def\period{3.4}     % cm per turn
  \def\xmax{8.0}
  \draw[engblue, thick, smooth, samples=120, domain=0:\xmax]
    plot (\x, {\helixA*sin(2*pi*\x/\period r)});
  \draw[engamber, thick, smooth, samples=120, domain=0:\xmax]
    plot (\x, {\helixA*sin(2*pi*\x/\period r + pi r)});
  % Base-pair rungs every 0.34 cm (= 0.34 nm / bp at 1 nm = 1 cm)
  \foreach \k in {0,1,...,23} {%
    \pgfmathsetmacro\xpos{0.17 + 0.34*\k}
    \pgfmathsetmacro\ya{\helixA*sin(2*pi*\xpos/\period r)}
    \pgfmathsetmacro\yb{\helixA*sin(2*pi*\xpos/\period r + pi r)}
    \draw[enggray, very thin] (\xpos,\ya) -- (\xpos,\yb);
  }
  % Annotations
  \draw[<->, thick] (0, -1.25) -- (3.4, -1.25)
    node[midway, fill=white, inner sep=1pt, font=\small, align=center]
    {pitch $3.4$ nm \\ ($10$ bp / turn)};
  \draw[<->, thick] (3.4, -1.25) -- (3.74, -1.25)
    node[below=2pt, font=\footnotesize, align=center] {rise \\ $0.34$ nm/bp};
  \draw[<->, thick] (-0.4, -\helixA) -- (-0.4, \helixA)
    node[midway, fill=white, inner sep=1pt, font=\small, rotate=90, align=center]
    {diameter $2$ nm};
  % Major / minor groove indication
  \node[engred, font=\footnotesize] at (1.7, 1.35) {major groove ($2.2$ nm wide)};
  \node[engred, font=\footnotesize] at (3.4, -0.05) {minor groove ($1.2$ nm wide)};
  % Strand polarity labels
  \node[engblue, anchor=west, font=\footnotesize] at (\xmax+0.05, 0.6) {$5' \to 3'$};
  \node[engamber, anchor=west, font=\footnotesize] at (\xmax+0.05, -0.6) {$3' \to 5'$};
\end{tikzpicture}

\caption{B-form DNA double helix. Two antiparallel sugar-phosphate
strands wrap around a common axis with a rise of $0.34$ nm per base
pair, ten base pairs per turn, an outer diameter of $2$ nm, and
major and minor grooves of $2.2$ nm and $1.2$ nm width respectively.
The grooves are the surfaces along which DNA-binding proteins read
the sequence.}
\label{fig:dna-geometry}
\end{figure}

The information content per base pair is $\log_2 4 = 2\,\text{bits}$.
The human genome is $3.2 \times 10^9\,\text{base pairs}$ in each haploid
set, or $6.4 \times 10^9\,\text{bits} = 0.8\,\text{GB}$ per diploid
nucleus. Compressed by removing repetitive sequence, the irreducible
information content is closer to $400$ to $700\,\text{MB}$. The
\emph{E.~coli} genome is $4.6 \times 10^6$ base pairs, or
$1.15\,\text{MB}$. The bacteriophage $\lambda$ has $48\,\text{kbp}$
($12\,\text{kB}$); a small viroid (a non-coding plant RNA) has
$250$ to $400\,\text{bases}$ ($60$ to $100\,\text{bytes}$). The
information densities cover seven orders of magnitude across the
biosphere.

\engterm{Empirical entropy of real genomes}. The $2$-bit upper bound
assumes uniform i.i.d.\ nucleotides. Real genomes deviate from this
limit in two ways. The single-nucleotide composition is GC-skewed
(human nuclear DNA is approximately $41\%$ G+C; \emph{E. coli} K-12
is $50.8\%$; the malaria parasite \emph{Plasmodium falciparum} is
$19.4\%$), which reduces the per-base entropy
$H_1 = -\sum_i p_i \log_2 p_i$ below $2$ bits. For human DNA the
entropy reduction is small ($H_1 \approx 1.99$ bits/bp); for
\emph{P.~falciparum} the AT bias drops it to $H_1 \approx 1.43$.
The conditional entropy $H_{2 \mid 1}$ accounting for neighbour
correlations falls further because of CpG suppression: in mammalian
genomes the CpG dinucleotide occurs at approximately $20\%$ of its
expected frequency, the consequence of cytosine methylation followed
by spontaneous deamination to thymine. The chapter's
\texttt{dna\_information\_content.py} script computes
$H_1$ and $H_{2 \mid 1}$ for an input FASTA and reports the
implied storage requirement.

\engterm{The molecular engineering of the helix}. The double helix is
mechanically stiff at lengths below its persistence length
$L_p \approx 50\,\text{nm}$ ($\approx 150\,\text{bp}$) and flexible
at longer lengths. The persistence length sets the geometric
limitations on how the genome can fold and how it can be threaded
through proteins. A $3.2 \times 10^9\,\text{bp}$ human genome
stretched end to end would be $1\,\text{m}$ long; the cell packages it
into a $5\,\si{\micro\meter}$ nucleus by wrapping the DNA around
histone proteins to form nucleosomes (each nucleosome is a
$\sim 11\,\text{nm}$ disc wrapping $147\,\text{bp}$), assembling
nucleosomes into a $30\,\text{nm}$ fibre, and folding the fibre into
higher-order loops and domains. The packaging factor is
$\sim 10^4$ in interphase chromatin and approximately $10^5$ in
metaphase chromosomes. The packaging is dynamic: the cell unwinds
regions of chromatin to allow access to transcription and
replication machinery, and re-winds them when access is no longer
needed. Chromatin structure is a control variable in gene expression
that Section~3.3 returns to.

\begin{table}[h]
\centering
\caption{Information-storage densities across the biosphere. The
``raw'' column reports $2 \log_2$ of the genome size in base pairs;
the ``irreducible'' column estimates the compressed content after
removing repetitive sequence and accounting for empirical entropy.
The biological design space spans seven orders of magnitude.}
\label{tab:genome-sizes}
\small
\begin{tabular}{lrrl}
\toprule
Organism & Genome (bp) & Raw (bytes) & Notes \\
\midrule
Viroid (PSTVd)            & $3.6 \times 10^2$ & $90$           & smallest known infectious nucleic acid \\
Bacteriophage $\phi$X174  & $5.4 \times 10^3$ & $1.4\,\text{kB}$ & first DNA genome sequenced (Sanger, 1977) \\
\emph{Mycoplasma genitalium} & $5.8 \times 10^5$ & $145\,\text{kB}$ & smallest known free-living organism \\
\emph{Escherichia coli} K-12 & $4.6 \times 10^6$ & $1.15\,\text{MB}$ & $\sim 4288$ protein-coding genes \\
\emph{Saccharomyces cerevisiae} & $1.2 \times 10^7$ & $3.0\,\text{MB}$ & first eukaryote sequenced (1996) \\
\emph{Drosophila melanogaster} & $1.4 \times 10^8$ & $36\,\text{MB}$ & $\sim 13\,000$ genes \\
\emph{Homo sapiens}        & $3.2 \times 10^9$ & $800\,\text{MB}$ & $\sim 20\,000$ protein-coding genes \\
\emph{Paris japonica} (a flowering plant) & $1.5 \times 10^{11}$ & $38\,\text{GB}$ & largest known plant genome \\
\bottomrule
\end{tabular}
\end{table}

\engterm{Why DNA, not RNA}. The cell uses DNA for stable information
storage and RNA for working transcripts. The choice is not arbitrary:
the absence of the $2'$-hydroxyl in deoxyribose makes DNA
substantially more chemically stable than RNA (the $2'$-OH catalyses
backbone hydrolysis in RNA), and the use of thymine rather than
uracil allows the cell to distinguish original-sequence thymines from
cytosine deamination products (cytosine deaminates spontaneously to
uracil at a measurable rate; if the genome used uracil, the deaminations
would be invisible to repair machinery). The substitution of thymine
for uracil is one of the cell's most elegant engineering choices: a
single methyl group provides the error-detection signature that
makes large-genome heredity feasible.

\engterm{Replication}. The double helix is replicated by separating
the two strands and using each as a template for synthesis of a new
complementary strand. The reaction is catalysed by DNA polymerases
that add nucleotides to a growing $3'$ end at a rate of approximately
$1000\,\text{nucleotides per second}$ in bacterial replication and
$100$ to $1000$ in eukaryotic replication. Bacterial replication is
bidirectional from a single origin (oriC in \emph{E.~coli}); eukaryotic
replication is bidirectional from thousands of origins distributed
along each chromosome. The leading strand is synthesised continuously;
the lagging strand is synthesised discontinuously in short
\engterm{Okazaki fragments} that are then joined by ligase. The
machinery (helicase, topoisomerase, primase, polymerase, sliding clamp,
ligase, and a dozen other accessory proteins) assembles into a
replication fork that resembles, in its engineering, a coordinated
multi-tool assembly machine. The Meselson-Stahl experiment
\cite{paper:v6c03-meselson-stahl-1958} established that the mechanism
is semiconservative (each daughter molecule carries one old and one
new strand) by density-gradient separation of nitrogen-labelled DNA;
the experiment is a canonical example of inferring molecular mechanism
from a single elegant measurement.

\engterm{Replication errors and proofreading}. DNA polymerase
incorporates the wrong nucleotide approximately once per $10^5$
incorporations from base-pairing fidelity alone. The polymerase's
$3' \rightarrow 5'$ exonuclease domain removes most of these errors
by proofreading immediately after incorporation, reducing the error
rate to approximately $10^{-7}$ per nucleotide. Mismatch-repair
machinery downstream catches most of the remaining errors, reducing
the error rate at fixation to approximately $10^{-9}$ to $10^{-10}$
per nucleotide per generation for fast-growing bacteria
\cite[ch.~25]{text:lehninger-biochemistry}. The cascade (base
selection, proofreading, mismatch repair) is an example of a
multi-stage error-correction architecture whose effective error rate
is the product of the per-stage error rates. Engineering parallels
in computer storage (parity, ECC, RAID) follow the same architectural
pattern. Figure~\ref{fig:error-cascade} traces the three stages and
their composition.

\begin{figure}[h]
\centering
% Multi-stage error-correction cascade for DNA replication.
\begin{tikzpicture}[
    >=Stealth,
    box/.style={draw, thick, rounded corners=2pt, minimum width=2.8cm,
      minimum height=1.0cm, align=center, font=\footnotesize, fill=engblue!10},
    rate/.style={draw, thick, fill=englime!15, minimum width=2.4cm,
      minimum height=0.7cm, font=\scriptsize, align=center},
    arrow/.style={->, thick},
  ]
  \node[box] (sel) {Base selection \\ \tiny Watson--Crick pairing};
  \node[rate, below=0.4cm of sel] (r1) {$\epsilon_1 \approx 10^{-5}$};
  \node[box, right=2cm of sel] (proof) {Proofreading \\ \tiny $3' \to 5'$ exonuclease};
  \node[rate, below=0.4cm of proof] (r2) {$\epsilon_2 \approx 10^{-2}$ \\ (residual after stage)};
  \node[box, right=2cm of proof] (mmr) {Mismatch repair \\ \tiny MutS / MutL / exo};
  \node[rate, below=0.4cm of mmr] (r3) {$\epsilon_3 \approx 10^{-2}$ \\ (residual after stage)};
  \node[box, right=2cm of mmr] (fix) {Fixation rate \\ \tiny per nucleotide};
  \node[rate, below=0.4cm of fix] (r4) {$\epsilon \approx 10^{-9}$};

  \draw[arrow] (sel) -- (proof);
  \draw[arrow] (proof) -- (mmr);
  \draw[arrow] (mmr) -- (fix);

  % Composition equation
  \node[font=\footnotesize, align=center] at (8.0, -2.3)
    {$\epsilon = \epsilon_1 \cdot \epsilon_2 \cdot \epsilon_3 \approx
      10^{-5} \cdot 10^{-2} \cdot 10^{-2} = 10^{-9}$};

  % Failure annotations
  \node[font=\scriptsize, engred, align=center, below=1.4cm of r1]
    {polymerase $\epsilon$/$\delta$ \\ mutations $\to$ ultramutators};
  \node[font=\scriptsize, engred, align=center, below=1.4cm of r2]
    {3'-exo loss \\ $\to$ POLE-mutated tumours};
  \node[font=\scriptsize, engred, align=center, below=1.4cm of r3]
    {MMR loss \\ $\to$ Lynch syndrome, MSI};
  \node[font=\scriptsize, engred, align=center, below=1.4cm of r4]
    {observed in fast-growing \\ bacteria and yeast};
\end{tikzpicture}

\caption{Multi-stage error-correction cascade for DNA replication.
Each stage reduces the residual error rate by a factor determined
by its molecular fidelity. The fixation rate
$\epsilon = \epsilon_1 \epsilon_2 \epsilon_3 \approx 10^{-9}$ matches
direct mutation-accumulation measurements in fast-growing bacteria.
Loss of any single stage (POLE 3'-exonuclease, MMR, polymerase
fidelity) produces a measurable mutator phenotype with the
characteristic mutational signature noted under each panel.}
\label{fig:error-cascade}
\end{figure}

\engterm{Worked example: error-rate composition}. Suppose a polymerase
has intrinsic base-selection error $\epsilon_1 = 3 \times 10^{-5}$,
the $3' \to 5'$ proofreading domain reduces the residual to
$\epsilon_1 \cdot \epsilon_2$ with $\epsilon_2 = 1/200$, and the
mismatch-repair machinery further reduces to
$\epsilon_1 \cdot \epsilon_2 \cdot \epsilon_3$ with
$\epsilon_3 = 1/500$. The composite fixation rate is
$\epsilon = 3 \times 10^{-5} \cdot 5 \times 10^{-3} \cdot 2 \times 10^{-3}
= 3 \times 10^{-10}$ per base. A $4.6 \times 10^6\,\text{bp}$
\emph{E.~coli} genome accumulates $1.4 \times 10^{-3}$ mutations per
generation, or one mutation every $\sim 720$ generations. The
order-of-magnitude agreement with measured mutation accumulation
\cite{paper:v6c03-lee-2012} is the engineering check on the cascade
model.

\begin{principle}[Genetic information is digital and error-corrected]
The cell stores information in a four-symbol digital code with a
multi-stage error-correction architecture whose effective per-base
error rate ($10^{-9}$ to $10^{-10}$) is several orders of magnitude
smaller than the per-stage error rate of any single step. The
architecture is the engineering condition that makes genomes of
$10^9$ to $10^{11}$ base pairs faithfully heritable across
generations. A reduction of the error-correction reliability by a
single factor of ten (loss of one proofreading or mismatch-repair
stage) produces measurable disease (the \engterm{mutator phenotype}
in cancer) within a few cell generations.
\end{principle}

\section{The central dogma engineering-first}\pathtag{core}

The \engterm{central dogma}, due to Francis Crick in 1957, states
that information in a cell flows in one direction: DNA $\rightarrow$
RNA $\rightarrow$ protein \cite{paper:v6c03-crick-1958}. The dogma has
well-known refinements (reverse transcription from RNA to DNA in
retroviruses; replication of RNA viruses) but the engineering reading
remains useful: the cell stores its information as DNA, transcribes it
to RNA when needed, and translates the RNA into the protein that
performs the cellular work. Figure~\ref{fig:central-dogma} summarises
the flow with its measured kinetic rates and per-step error rates.

\begin{figure}[h]
\centering
% Central-dogma flow diagram with measured kinetic rates.
\begin{tikzpicture}[
    >=Stealth,
    node distance=2.6cm,
    box/.style={draw, thick, rounded corners=2pt, minimum width=2.3cm,
      minimum height=1.1cm, align=center, font=\small},
    arrow/.style={->, thick, engblue},
    rev/.style={->, thick, dashed, engred},
    label/.style={font=\footnotesize, fill=white, inner sep=1pt, align=center},
  ]
  \node[box, fill=engblue!10] (dna)  {DNA \\ \footnotesize $2$ bits/bp};
  \node[box, fill=englime!20, right=of dna] (rna) {pre-mRNA \\ \footnotesize $5'$ cap, $3'$ poly(A)};
  \node[box, fill=englime!40, right=of rna] (mrna) {mature \\ mRNA};
  \node[box, fill=engamber!25, right=of mrna] (prot) {protein \\ \footnotesize folded};

  \draw[arrow] (dna) -- (rna)
    node[label, midway, above] {transcription \\ $\sim 50$ nt/s} ;
  \draw[arrow] (rna) -- (mrna)
    node[label, midway, above] {splicing \\ cap, tail};
  \draw[arrow] (mrna) -- (prot)
    node[label, midway, above] {translation \\ $\sim 20$ aa/s};

  % Replication self-loop on DNA
  \draw[arrow] (dna.north west) .. controls +(-0.9,0.9) and +(-0.9,-0.9) ..
    (dna.south west) node[label, midway, left] {replication \\ $\sim 1000$ nt/s};

  % Reverse transcription
  \draw[rev] (rna.south) to[bend right=20]
    node[label, midway, below] {reverse transcription \\ (retroviruses)} (dna.south);

  % Error-rate annotations under each step
  \node[font=\scriptsize, engred] at (dna.south) [yshift=-1.5cm] {err $10^{-9}$/bp};
  \node[font=\scriptsize, engred] at (rna.south) [yshift=-0.9cm] {err $10^{-5}$/nt};
  \node[font=\scriptsize, engred] at (prot.south) [yshift=-0.9cm] {err $10^{-3}$/codon};
\end{tikzpicture}

\caption{The central dogma with measured rates and error budgets.
Replication runs at $\sim 1000$ nt/s in \emph{E.~coli}; transcription
elongates at $\sim 50$ nt/s; translation elongates at $\sim 20$ aa/s
in bacteria. Error budgets compound across stages: replication is
heritable and so is policed to $10^{-9}$/bp; transcription is
transient (a faulty mRNA is replaced) and runs at $\sim 10^{-5}$/nt;
translation is even more transient and runs at $\sim 10^{-3}$/codon.
Reverse transcription (dashed) is the dogma's principal exception.}
\label{fig:central-dogma}
\end{figure}

\engterm{Transcription}. The cell copies a specific stretch of DNA
into a complementary RNA strand by means of RNA polymerase. The
polymerase binds to a \engterm{promoter} sequence upstream of the
gene, unwinds the DNA locally, and synthesises RNA at a rate of
approximately $50$ to $100\,\text{nucleotides per second}$. The RNA
strand is identical in sequence to the non-template DNA strand, with
uracil replacing thymine. Transcription terminates at a
\engterm{terminator} sequence (in bacteria) or by polyadenylation
signal (in eukaryotes). Bacterial and eukaryotic transcription differ
in detail (bacteria use a single RNA polymerase with sigma factors
for promoter recognition; eukaryotes use three distinct polymerases
plus a complex of general transcription factors at each promoter) but
share the engineering structure: a regulated synthesis of a
sequence-defined RNA copy from a sequence-defined DNA template.

\engterm{RNA processing}. In eukaryotes, the primary RNA transcript
undergoes three principal modifications before translation. A
$5'$ cap (a modified guanosine) is added to the front; a poly(A) tail
of $\sim 200$ adenines is added to the back; and \engterm{introns}
(non-coding sequences between coding \engterm{exons}) are removed by
the spliceosome. The processed mRNA is exported from the nucleus to
the cytoplasm for translation. Alternative splicing (different
combinations of exons included or excluded) allows a single gene to
encode multiple protein variants; the human genome's
$\sim 20\,000$ protein-coding genes encode an estimated $\sim 100\,000$
distinct protein variants through alternative splicing.

\engterm{Translation}. The mRNA is decoded by ribosomes into a
protein. The ribosome is a large ribonucleoprotein complex
($\sim 2.5\,\text{MDa}$ in bacteria, $\sim 4\,\text{MDa}$ in
eukaryotes) consisting of two subunits with three transfer-RNA
(tRNA) binding sites. Each tRNA carries an amino acid at one end and
a three-nucleotide anticodon at the other; the anticodon base-pairs
with a complementary codon in the mRNA. The ribosome moves along the
mRNA in steps of three nucleotides, with each step adding one amino
acid to the growing peptide chain. Translation proceeds at
$15$ to $20$ amino acids per second in bacteria, slightly slower in
eukaryotes. The ribosome reads the mRNA from the $5'$ end to the
$3'$ end, starts at an AUG start codon, and stops at one of three
stop codons (UAA, UAG, UGA). The completed peptide is released and
folds into a functional protein (Chapter~\ref{vol06:ch04}).

\engterm{Error rates}. Translation makes errors at a rate of
approximately $10^{-3}$ to $10^{-4}$ per codon, dominated by tRNA
mischarging (the aminoacyl-tRNA synthetase loading the wrong amino
acid onto a tRNA) and by codon-anticodon mispairing at the ribosome.
The rate is much higher than the DNA replication rate because the
errors are transient (one mistranslated protein molecule among many
correctly translated copies) and the cost of correcting them at
synthesis time exceeds the cost of tolerating them in the protein
pool. The cell's error-rate budget at each information-processing
step is engineered to the cost of failure at that step: replication
errors are heritable and so are corrected aggressively; translation
errors are transient and so are tolerated more.

\engterm{Kinetic model of mRNA and protein abundance}. The two-stage
flow admits a coupled linear ODE model with measured constants:
\[
\frac{dm}{dt} = k_\text{tx} - \gamma_m\, m,
\qquad
\frac{dp}{dt} = k_\text{tl}\, m - \gamma_p\, p.
\]
For an \emph{E. coli} gene with $k_\text{tx} = 0.1$ mRNA/s,
$k_\text{tl} = 0.05$ protein/(mRNA $\cdot$ s), mRNA half-life
$\tau_{1/2}^m = 5$ minutes (so $\gamma_m = \ln 2/300\,\text{s}$), and
protein half-life $\tau_{1/2}^p = 20$ hours (so
$\gamma_p = \ln 2/72\,000\,\text{s}$), the steady-state copy numbers
are $m_\text{ss} = k_\text{tx}/\gamma_m \approx 43$ mRNA and
$p_\text{ss} = k_\text{tl} m_\text{ss}/\gamma_p \approx 220\,000$
protein. The mRNA time constant ($\sim 7$ min) sets the speed of
transcriptional response; the protein time constant ($\sim 29$ h) sets
the time over which an environmental change produces a measurable
shift in protein abundance. The chapter's
\texttt{central\_dogma\_kinetics.py} script integrates the system
and supports a step-induction option, illustrating that the protein
response to a transcription burst is delayed by the protein half-life
not the mRNA half-life. The mismatch in time constants between mRNA
and protein is a control-engineering trade-off: a short-lived mRNA
allows rapid retuning of gene expression at low protein cost; a
long-lived protein amortises the synthesis cost over many cell-state
changes.

\section{Gene expression and regulation}\pathtag{standard}

A cell's genome contains the parts list of every protein the cell
might ever need. The cell does not synthesise all of them at once;
different cell types, different developmental stages, and different
environmental conditions require different protein complements. The
machinery that selects which genes are expressed and at what level is
the \engterm{gene-expression regulatory network}.

\engterm{Transcriptional regulation}. The dominant level of control is
the rate of transcription initiation at each gene. The control is
exerted by \engterm{transcription factors}, proteins that bind
specific DNA sequences near a gene's promoter and either recruit or
exclude the RNA polymerase. A typical bacterial gene's promoter is
regulated by one to a few transcription factors; a typical
eukaryotic gene's promoter is regulated by a dozen or more
transcription factors plus distal enhancer elements that may be tens
of kilobases away. The combinatorial logic of multiple transcription
factors at a single promoter implements logical operations (AND, OR,
NOT) on the regulatory inputs; the network architecture computes
gene expression as a function of the cell's state.
Figure~\ref{fig:regulation-circuit} shows the principal architectural
elements.

\begin{figure}[h]
\centering
% Gene-regulation circuit: enhancer, repressor, activator, TATA, Pol II
\begin{tikzpicture}[>=Stealth, scale=1.0,
    tf/.style={draw, thick, rounded corners=3pt, minimum width=1.4cm, minimum height=0.7cm,
      align=center, font=\footnotesize},
    site/.style={draw, fill=enggray!20, minimum width=0.9cm, minimum height=0.4cm,
      font=\scriptsize, inner sep=1pt},
  ]
  % DNA double line
  \draw[engblue, very thick] (-0.5, 0) -- (12, 0);
  \draw[engblue, very thick] (-0.5, -0.25) -- (12, -0.25);

  % Cis-regulatory elements along DNA
  \node[site] (enh) at (0.6, -0.125) {ENH};
  \node[site] (rep) at (3.0, -0.125) {OP};
  \node[site] (uas) at (5.4, -0.125) {UAS};
  \node[site] (tata) at (7.6, -0.125) {TATA};
  \node[site, fill=englime!30] (tss) at (8.7, -0.125) {+1};
  \node[draw, fill=engamber!20, minimum width=2.5cm, minimum height=0.4cm,
        font=\scriptsize] (gene) at (10.5, -0.125) {coding sequence};

  % Trans-factors above DNA
  \node[tf, fill=englime!30, above=0.7cm of enh] (act) {activator};
  \node[tf, fill=engred!30,  above=0.7cm of rep] (reps) {repressor};
  \node[tf, fill=englime!30, above=0.7cm of uas] (act2) {coactivator};
  \node[tf, fill=engblue!20, above=1.3cm of tata, minimum width=2.0cm] (poly)
    {RNA Pol II \\ + GTFs};

  \draw[->, thick, englime] (act.south) -- (enh.north);
  \draw[->, thick, engred]  (reps.south) -- (rep.north);
  \draw[->, thick, englime] (act2.south) -- (uas.north);
  \draw[->, thick, engblue] (poly.south) -- (tata.north);

  % Long-range looping (enhancer-promoter)
  \draw[engamber, thick, dashed, ->] (enh.north west) .. controls +(-0.5,2.2) and +(-2,1.2) ..
    (poly.west) node[midway, above, font=\scriptsize, engamber] {DNA looping};

  % Outputs
  \draw[->, very thick, englime] (tss.south) -- ++(0,-0.7)
    node[anchor=north, font=\footnotesize] {ON (activator dominant)};
  \draw[->, very thick, engred] (rep.south) -- ++(0,-1.4)
    node[anchor=north, font=\footnotesize] {OFF (repressor dominant)};
\end{tikzpicture}

\caption{Eukaryotic-promoter regulatory architecture. An enhancer
(ENH) and an upstream activating sequence (UAS) bind activators that
contact the polymerase complex through DNA looping; an operator (OP)
binds a repressor that blocks the polymerase. The TATA box (or
equivalent core-promoter element) and the transcription start site
($+1$) locate the polymerase. The expressed-versus-silent state of
the gene is the integrated logic of these binding events;
\emph{E.~coli} promoters generally have a simpler one-to-two
transcription-factor architecture.}
\label{fig:regulation-circuit}
\end{figure}

\engterm{The lac operon}. The canonical bacterial example is the
\emph{lac} operon of \emph{E. coli}, characterised by Jacob and
Monod \cite{paper:v6c03-jacob-monod-1961}. Three genes encoding
$\beta$-galactosidase (\emph{lacZ}), lactose permease (\emph{lacY}),
and a thiogalactoside transacetylase (\emph{lacA}) are co-transcribed
from a single promoter under the control of two regulators: the LacI
repressor binds the operator and blocks transcription in the absence
of lactose, and the CAP-cAMP complex activates transcription only
when glucose is unavailable. The output is a Boolean function of two
inputs: ON only when lactose is present \emph{and} glucose is absent.
The architecture is one of the earliest examples of a regulated
circuit whose engineering function was reverse-engineered from
biochemistry, and it remains the canonical lecture-room example of
combinatorial transcriptional logic.

\engterm{Network architecture}. The cell's transcription network has
been mapped exhaustively for \emph{E.~coli} (approximately $300$
known transcription factors regulating $4288$ genes) and for several
other model organisms. The network is sparse: each transcription
factor regulates a few to a few hundred targets; each gene is
regulated by a few transcription factors. The network contains
recurring motifs whose engineering function has been worked out:
auto-regulation (a transcription factor regulating its own gene,
either negatively for noise reduction or positively for switch-like
behaviour), the feed-forward loop (Section~\ref{vol06:ch01} introduced
the type-1 coherent loop, with its delay and noise-filtering behaviour),
the single-input module (a master regulator controlling a set of
co-regulated targets), and the dense overlapping regulon (multiple
factors jointly controlling overlapping target sets)
\cite[ch.~6]{text:alon-systems-biology}.

\engterm{Post-transcriptional regulation}. The cell controls protein
abundance not only by transcription rate but also by mRNA stability
(half-lives ranging from minutes for short-lived mRNAs to days for
stable structural mRNAs), translation efficiency (set by the mRNA's
$5'$ untranslated region and by sequence-specific binding factors),
protein stability (set by the protein's degradation signals,
including ubiquitin tagging for proteasomal degradation), and
post-translational modification (phosphorylation, acetylation,
methylation). The cell's effective gene-expression control surface
combines all these levels, with the result that the steady-state
abundance of a protein can vary by factors of $10^4$ to $10^6$ across
cell states without changing the underlying DNA sequence.

\engterm{Epigenetic regulation}. Some changes in gene expression
persist across cell divisions without changes in DNA sequence. The
two principal mechanisms are DNA methylation (covalent attachment of
methyl groups to cytosines, suppressing gene expression in mammals)
and histone modification (acetylation, methylation, and other
covalent marks on the histone proteins around which DNA is wrapped,
affecting chromatin accessibility). Together these constitute the
\engterm{epigenome}, a layer of regulatory information that is
copied along with the DNA during replication and that maintains
cell-type-specific gene expression patterns. Epigenetic engineering
is an active area of biomanufacturing and therapy at the time of
writing (current as of 2025).

\engterm{Worked example: induction time of a controlled gene}. A
strain engineered to express green fluorescent protein (GFP) under
the \emph{lac} promoter is induced by adding IPTG to the medium. The
transcript begins to accumulate immediately; the protein lags. With
a representative GFP mRNA half-life of $5$ minutes and a GFP protein
half-life of $26$ hours in \emph{E. coli}, the protein reaches
$50\%$ of steady state at
$t_{1/2} = \tau_p \ln 2 \approx 18\,\text{h}$, where
$\tau_p = 1/\gamma_p$. The mismatch between the laboratory perception
of induction (the bottle of cells glows green within an hour) and the
nominal steady state (eighteen hours of accumulation later) is the
source of repeated confusion in biotech process control. The protein
signal at any earlier time is set by the integration
$p(t) = p_\text{ss}(1 - e^{-\gamma_p t})$ once the mRNA pool has
equilibrated (after a few mRNA half-lives), an order-of-magnitude
calculation that the chapter's kinetics script reproduces directly.

\section{The genetic code}\pathtag{standard}

The \engterm{genetic code} is the mapping from nucleotide triplets
(codons) to amino acids and stop signals. There are $4^3 = 64$
possible codons. Three are stop codons (UAA, UAG, UGA); the remaining
$61$ are sense codons that specify the $20$ standard amino acids.
The code is therefore degenerate: most amino acids are specified by
two to six different codons, and only methionine (AUG, the standard
start codon) and tryptophan (UGG) are specified by a single codon.
The complete code is shown in Figure~\ref{fig:genetic-code}.

\begin{figure}[h]
\centering
% Standard genetic code as a 4x4 colour-coded grid, second-position column,
% first-position row, third-position label inside each cell.
\begin{tikzpicture}[
    cell/.style={draw, minimum width=2.0cm, minimum height=1.2cm, align=center,
      font=\footnotesize, inner sep=2pt},
    header/.style={font=\bfseries\small, align=center},
    % colour scheme: hydrophobic (yellow), polar (cyan), basic (blue),
    % acidic (red), special (gray), stop (black)
    hphobic/.style={cell, fill=engamber!20},
    polar/.style={cell, fill=englime!25},
    basic/.style={cell, fill=engblue!25},
    acidic/.style={cell, fill=engred!25},
    special/.style={cell, fill=enggray!25},
    stop/.style={cell, fill=engblack!70, text=white},
  ]
  % Column headers (second position)
  \node[header] at (1, 5.3) {U};
  \node[header] at (3, 5.3) {C};
  \node[header] at (5, 5.3) {A};
  \node[header] at (7, 5.3) {G};
  % Row headers (first position)
  \node[header] at (-0.6, 4.4) {U};
  \node[header] at (-0.6, 3.0) {C};
  \node[header] at (-0.6, 1.6) {A};
  \node[header] at (-0.6, 0.2) {G};

  % --- First position U ---
  \node[hphobic] at (1, 4.4) {Phe UU\textbf{U,C} \\ Leu UU\textbf{A,G}};
  \node[hphobic] at (3, 4.4) {Ser UC\textbf{N}};
  \node[polar]   at (5, 4.4) {Tyr UA\textbf{U,C} \\ STOP UA\textbf{A,G}};
  \node[special] at (7, 4.4) {Cys UG\textbf{U,C} \\ STOP UGA \\ Trp UGG};
  % --- First position C ---
  \node[hphobic] at (1, 3.0) {Leu CU\textbf{N}};
  \node[hphobic] at (3, 3.0) {Pro CC\textbf{N}};
  \node[basic]   at (5, 3.0) {His CA\textbf{U,C} \\ Gln CA\textbf{A,G}};
  \node[basic]   at (7, 3.0) {Arg CG\textbf{N}};
  % --- First position A ---
  \node[hphobic] at (1, 1.6) {Ile AU\textbf{U,C,A} \\ Met AUG (start)};
  \node[polar]   at (3, 1.6) {Thr AC\textbf{N}};
  \node[polar]   at (5, 1.6) {Asn AA\textbf{U,C} \\ Lys AA\textbf{A,G}};
  \node[polar]   at (7, 1.6) {Ser AG\textbf{U,C} \\ Arg AG\textbf{A,G}};
  % --- First position G ---
  \node[hphobic] at (1, 0.2) {Val GU\textbf{N}};
  \node[polar]   at (3, 0.2) {Ala GC\textbf{N}};
  \node[acidic]  at (5, 0.2) {Asp GA\textbf{U,C} \\ Glu GA\textbf{A,G}};
  \node[polar]   at (7, 0.2) {Gly GG\textbf{N}};

  % Legend
  \node[font=\footnotesize, anchor=west] at (-1.0, -1.0)
    {Legend: \tikz\node[cell, minimum width=0.4cm, minimum height=0.3cm,
       fill=engamber!20]{};~hydrophobic ~~
    \tikz\node[cell, minimum width=0.4cm, minimum height=0.3cm,
       fill=englime!25]{};~polar/uncharged ~~
    \tikz\node[cell, minimum width=0.4cm, minimum height=0.3cm,
       fill=engblue!25]{};~basic ~~
    \tikz\node[cell, minimum width=0.4cm, minimum height=0.3cm,
       fill=engred!25]{};~acidic ~~
    \tikz\node[cell, minimum width=0.4cm, minimum height=0.3cm,
       fill=engblack!70]{};~stop};
\end{tikzpicture}

\caption{The standard nuclear genetic code. Cells are colour-coded
by the chemical class of the encoded amino acid: hydrophobic
(amber), polar uncharged (green), basic (blue), acidic (red), special
(grey), stop (black). The third codon position (the wobble position)
is the most permissive; single-nucleotide errors at the third
position are disproportionately likely to leave the amino acid
unchanged or to substitute a chemically similar residue. The full
table is available in \texttt{data/genetic-code.csv}.}
\label{fig:genetic-code}
\end{figure}

\engterm{Code structure}. The degeneracy follows a regular pattern.
The third codon position (the ``wobble position'') is the most
variable; for many amino acids, all four nucleotides at the third
position specify the same amino acid (the third-position degeneracy
is fourfold). The first and second positions are more constrained.
The codon table groups codons by physical and chemical properties of
the encoded amino acid: codons starting with U specify hydrophobic
amino acids (phenylalanine, leucine, isoleucine, methionine, valine);
codons with C in the second position specify small or polar amino
acids; codons with A in the second position specify polar or charged
amino acids. The pattern is not random: a single-nucleotide
substitution is statistically more likely to substitute a chemically
similar amino acid (a conservative substitution) than to substitute a
chemically dissimilar amino acid (a radical substitution). The bias is
the code's principal error-mitigation feature.

\engterm{Deciphering the code}. The code was decoded experimentally
between 1961 and 1966. Nirenberg and Matthaei \cite{paper:v6c03-nirenberg-matthaei-1961}
demonstrated in 1961 that a synthetic poly(U) RNA directs synthesis of
polyphenylalanine, identifying UUU as a codon for phenylalanine. The
remaining assignments were worked out by similar synthetic-RNA
experiments and by ribosome-binding assays over the following five
years. The five-year deciphering project is a model of how an
information-theoretic question (what is the codon table?) was answered
by a series of bench experiments whose interpretation required a
physical model of the ribosome and the tRNA pool.

\engterm{Universality}. The code is almost universal: the same codon
specifies the same amino acid in every well-studied organism on Earth.
The few known exceptions (variant codes in mitochondria, in some
ciliates, in some bacterial lineages) involve a small number of
codons. The universality reflects the evolutionary cost of changing
the code: once a code is in place, any change requires simultaneous
changes in the tRNA pool, the aminoacyl-tRNA synthetases, and every
gene whose codons were affected. The code is therefore an example of
\engterm{frozen accident}: a feature that was variable at the origin
of life but has been frozen since.

\engterm{Designed by evolution}. The genetic code's structure (the
specific assignment of codons to amino acids) is the result of
evolutionary optimisation against several competing pressures:
translation accuracy (similar codons should specify similar amino
acids, so that translation errors are tolerable), mutational
robustness (single-nucleotide mutations should preferentially produce
chemically similar substitutions), and biosynthetic cost (frequently
used amino acids should have codons that minimise the cell's resource
demand on tRNA synthesis). Freeland and Hurst's 1998 analysis
\cite{paper:v6c03-freeland-hurst-1998} compared the standard code to
$10^6$ random reassignments of the codon table to amino acids; the
standard code lies in the top $10^{-6}$ for minimising the average
chemical-property change under single-nucleotide mutation. The
reading is that the code was selected by evolution as an engineering
trade-off, not designed by deliberate optimisation. The Freeland-Hurst
result is the quantitative reading of the qualitative observation in
\cite[ch.~27]{text:lehninger-biochemistry} that the code is
``one in a million.''

\engterm{Codon optimisation in biomanufacturing}. When a gene is
expressed in a heterologous host (e.g.\ a human gene placed in
\emph{E. coli} for industrial protein production), the gene's codon
usage may not match the host's tRNA pool. Codons that are rare in the
host slow translation, cause ribosome stalling, and reduce yield.
The codon-adaptation index (CAI) of Sharp and Li
\cite{paper:v6c03-sharp-li-1987} quantifies the match between a gene's
codon usage and a reference set of highly expressed host genes; a CAI
near $1$ indicates good adaptation, a CAI near $0$ indicates poor
adaptation. Commercial gene-synthesis services routinely recode
heterologous genes to maximise CAI in the target host, often
increasing protein yield by factors of $2$ to $10$. The chapter's
\texttt{codon\_optimization.py} script computes CAI against an
\emph{E. coli} reference table and provides a greedy recoder.

\section{Mutation, repair, and error rates}\pathtag{core}

\engterm{Mutations}. A mutation is a change in the DNA sequence. The
classes:

\begin{itemize}[leftmargin=*]
\item \engterm{Point mutations} (single base changes), which divide
  further into transitions (purine $\leftrightarrow$ purine, or
  pyrimidine $\leftrightarrow$ pyrimidine) and transversions (purine
  $\leftrightarrow$ pyrimidine). Most spontaneous point mutations are
  transitions because the chemical mechanisms (cytosine deamination,
  guanine oxidation) preferentially produce transitions.
\item \engterm{Insertions and deletions} (indels) of one or more
  base pairs. Indels in coding regions cause frameshift mutations
  unless the indel length is a multiple of three; frameshifts
  typically destroy gene function.
\item \engterm{Large structural variants}: duplications, inversions,
  translocations, copy-number variants. These arise during recombination
  or by repair of double-strand breaks. They are individually rare but
  important contributors to genetic disease and to cancer.
\end{itemize}

\engterm{Mutation rates}. The spontaneous mutation rate per base pair
per generation is approximately $10^{-9}$ in bacteria, $10^{-10}$ in
unicellular eukaryotes, and $10^{-8}$ to $10^{-9}$ in mammalian
germline cells \cite[ch.~25]{text:lehninger-biochemistry}. Per genome
per generation, this gives roughly $0.005$ mutations per bacterial
generation, $0.003$ for yeast, and $\sim 100$ for the human germline.
The human-genome value translates to approximately $50$ to $70$ new
single-base variants per child compared to the parents
\cite{paper:v6c03-kong-2012}, a number that genomic sequencing now
confirms directly. Table~\ref{tab:mutation-rates} summarises rates
across twelve organisms; the data file
\texttt{data/mutation-rates.csv} carries the full table.

\begin{table}[h]
\centering
\caption{Per-base mutation rates for representative organisms (current
as of 2024). Rates are per base pair per generation (cellular
organisms) or per replication cycle (viruses). The four orders of
magnitude span between an RNA virus and a fast-growing bacterium
reflect the absence of effective proofreading in RNA-virus
polymerases.}
\label{tab:mutation-rates}
\small
\begin{tabular}{lrr}
\toprule
Organism & per-bp rate (per generation) & per-genome \\
\midrule
HIV-1 (per replication)       & $3.4 \times 10^{-5}$  & $0.33$ \\
SARS-CoV-2 (per replication)  & $1.0 \times 10^{-6}$  & $0.03$ \\
Bacteriophage $\lambda$       & $7.7 \times 10^{-8}$  & $3.7 \times 10^{-3}$ \\
\emph{Escherichia coli} K-12  & $2.2 \times 10^{-10}$ & $1.0 \times 10^{-3}$ \\
\emph{Saccharomyces cerevisiae} & $3.3 \times 10^{-10}$ & $4.0 \times 10^{-3}$ \\
\emph{Caenorhabditis elegans} & $2.7 \times 10^{-9}$  & $0.27$ \\
\emph{Drosophila melanogaster} & $2.8 \times 10^{-9}$ & $0.40$ \\
\emph{Mus musculus} (germline) & $5.4 \times 10^{-9}$ & $15$ \\
\emph{Homo sapiens} (germline) & $1.2 \times 10^{-8}$ & $70$ \\
\emph{H.~sapiens} (somatic, colon) & $4.3 \times 10^{-9}$ & $13$ per division \\
\bottomrule
\end{tabular}
\end{table}

\engterm{DNA damage and repair}. Beyond replication errors, the cell's
DNA is continuously damaged by chemical and physical agents: oxidation
of bases (the most common spontaneous lesion, with approximately
$10\,000$ oxidative lesions per human cell per day), deamination,
alkylation, and ultraviolet-induced pyrimidine dimers. The cell repairs
these lesions through several pathways: base excision repair (for
oxidised and small chemical lesions), nucleotide excision repair (for
bulky lesions like UV-induced dimers), mismatch repair (for replication
errors that escape proofreading), and double-strand-break repair by
homologous recombination or non-homologous end joining (for the
most severe lesions). Each pathway has its own engineering
characteristics (speed, fidelity, applicability), and the cell selects
the appropriate pathway based on the lesion type and the cell-cycle
phase. Figure~\ref{fig:repair-pathways} routes the four principal
lesion classes to the corresponding repair machinery and indicates the
pathological consequence of each pathway's failure.

\begin{figure}[h]
\centering
% Mutation + repair flowchart: lesion classes routed to BER / NER / MMR / NHEJ / HR.
\begin{tikzpicture}[
    >=Stealth, node distance=1.4cm and 2.1cm,
    lesion/.style={draw, thick, rounded corners=2pt, fill=engred!15,
      minimum width=2.4cm, minimum height=0.9cm, align=center, font=\footnotesize},
    repair/.style={draw, thick, rounded corners=2pt, fill=englime!25,
      minimum width=2.4cm, minimum height=0.9cm, align=center, font=\footnotesize},
    outcome/.style={draw, thick, rounded corners=2pt, fill=engblue!15,
      minimum width=2.4cm, minimum height=0.9cm, align=center, font=\footnotesize},
    arrow/.style={->, thick},
  ]
  % Lesion column
  \node[lesion] (mismatch) {replication \\ mismatch};
  \node[lesion, below=of mismatch] (oxbase) {oxidised / \\ alkylated base};
  \node[lesion, below=of oxbase] (uv) {UV pyrimidine \\ dimer / bulky adduct};
  \node[lesion, below=of uv] (dsb) {double-strand \\ break};

  % Repair column
  \node[repair, right=of mismatch] (mmr) {MMR \\ \tiny MutS, MutL, exo};
  \node[repair, right=of oxbase]   (ber) {BER \\ \tiny glycosylase, AP endo};
  \node[repair, right=of uv]       (ner) {NER \\ \tiny XPA-G, ERCC};
  \node[repair, right=of dsb]      (hr)  {HR (S/G2) \\ \tiny BRCA1/2, RAD51};
  \node[repair, below=0.5cm of hr] (nhej){NHEJ (G1) \\ \tiny Ku70/80, DNA-PK};

  % Outcomes
  \node[outcome, right=of mmr] (fix1) {restored \\ sequence};
  \node[outcome, right=of ber] (fix2) {restored \\ base};
  \node[outcome, right=of ner] (fix3) {patched \\ sequence};
  \node[outcome, right=of hr]  (fix4) {error-free \\ rejoin};
  \node[outcome, right=of nhej, fill=engamber!25] (fix5) {error-prone \\ rejoin};

  % Failures
  \node[draw, dashed, thick, fill=engred!8, font=\scriptsize, align=center,
        below=1.2cm of fix5, minimum width=10cm] (fail)
    {\textbf{Repair failure pathologies.}
     MMR loss $\to$ Lynch syndrome / MSI tumours \quad
     NER loss $\to$ xeroderma pigmentosum \quad
     HR loss $\to$ BRCA-driven cancers \quad
     NHEJ misjoin $\to$ translocations (BCR-ABL)};

  \draw[arrow] (mismatch) -- (mmr); \draw[arrow] (mmr) -- (fix1);
  \draw[arrow] (oxbase)   -- (ber); \draw[arrow] (ber) -- (fix2);
  \draw[arrow] (uv)       -- (ner); \draw[arrow] (ner) -- (fix3);
  \draw[arrow] (dsb.east) -- ++(0.5,0) |- (hr.west);
  \draw[arrow] (dsb.east) -- ++(0.5,0) |- (nhej.west);
  \draw[arrow] (hr) -- (fix4); \draw[arrow] (nhej) -- (fix5);

  \draw[arrow, dashed, engred] (fix5.south) -- (fail.north);
\end{tikzpicture}

\caption{Lesion-to-repair routing in the human cell. Mismatch repair
(MMR) corrects replication errors; base excision repair (BER) handles
oxidised and small chemical lesions; nucleotide excision repair (NER)
removes bulky adducts including UV pyrimidine dimers; homologous
recombination (HR) and non-homologous end joining (NHEJ) repair
double-strand breaks, with HR error-free in S/G2 and NHEJ
error-prone in G1. Each pathway has a characteristic disease
signature when it fails (MMR $\to$ Lynch syndrome, NER $\to$
xeroderma pigmentosum, HR $\to$ BRCA-driven cancers).}
\label{fig:repair-pathways}
\end{figure}

\engterm{Error rates in engineering context}. The cell's
$10^{-9}$ to $10^{-10}$ per-base error rate is comparable to the
bit-error rate of high-quality hard disk storage
($\sim 10^{-15}$ with ECC, but $\sim 10^{-6}$ without). The cell's
multi-stage architecture (replication accuracy, proofreading,
mismatch repair, damage repair) achieves its low overall error rate
by composition of imperfect stages, the same architecture pattern as
modern storage systems. The engineering reading is that DNA storage,
treated as an engineered system, is one of the most reliable
information-storage technologies ever produced, with retention times
of $10^4$ to $10^6$ years under appropriate conditions (the recovery
of ancient DNA from specimens up to $10^6$ years old demonstrates the
shelf life).

\engterm{Worked example: estimating the human mutation rate from a
trio}. A whole-genome sequence of a child plus both parents at
$30 \times$ coverage typically yields $\sim 70$ candidate de novo
variants after filtering for sequencing artefacts, mosaicism in the
parents, and recurrent false-positive sites. After validation by an
orthogonal method (typically Sanger sequencing or amplicon
re-sequencing on a different platform), approximately $50$ to $60$
true de novo variants remain. Dividing by the $\sim 2.8 \times 10^9$
callable bases in the diploid genome (allowing for $90\%$ callability)
gives an estimate of
$\sim 1.0\text{--}1.1 \times 10^{-8}$ per bp per generation. The
chapter's \texttt{mutation\_rate\_estimator.py} provides the
calculation, the Wilson $95\%$ confidence interval, and a
synthetic-trio simulator for validation. The result agrees with the
Iceland trio analysis of Kong et al. \cite{paper:v6c03-kong-2012},
which gives $1.20 \times 10^{-8}$ with paternal age accounting for
roughly two new mutations per additional paternal year.

\section{Modern sequencing technology}\pathtag{standard}

The ability to read DNA sequences has transformed biology and
biotechnology over the past two decades. The chapter surveys the
principal technologies in the order they have emerged.

\engterm{Sanger sequencing}. The \engterm{chain-termination method}
developed by Frederick Sanger in 1977 \cite{paper:v6c03-sanger-1977}
used DNA polymerase to synthesise a complementary strand in the
presence of low concentrations of fluorescently labelled
dideoxynucleotides. Each ddNTP terminates the chain upon
incorporation; the resulting mixture of variable-length fluorescently
labelled fragments is separated by electrophoresis and the sequence
read from the fragment lengths. Sanger sequencing achieves read
lengths of $500$ to $1000$ bases at error rates of $\sim 10^{-4}$,
and remained the dominant technology from $1977$ to approximately
$2007$. The Human Genome Project (1990 to 2003) was completed with
Sanger sequencing at a final cost of approximately $\$3$ billion for
the first reference genome.

\engterm{Short-read next-generation sequencing}. From 2005 onward,
Illumina's reversible-terminator chemistry (and earlier Solexa
technology) \cite{paper:v6c03-bentley-2008} drove sequencing costs
down by approximately six orders of magnitude. The platform sequences
millions of short ($50$ to $300\,\text{base}$) reads in parallel by
cycling fluorescent nucleotides through a flow cell of immobilised
DNA fragments. Error rates are approximately $10^{-3}$, mostly
substitutions; throughputs are $10^{12}$ to $10^{13}$ bases per run.
The cost per human genome fell from $\$3\,\text{billion}$ in 2003 to
$\$10\,000$ in 2010 to $\$1000$ in 2014 to under $\$200$ in 2023
(current as of 2025) \cite{paper:v6c03-nhgri-cost-2023}. The cost
trajectory has substantially outpaced Moore's law in the same period
\cite[ch.~12]{text:alon-systems-biology}. Figure~\ref{fig:seq-cost}
shows the NHGRI cost-per-genome series against a Moore's-law
reference.

\begin{figure}[h]
\centering
% Cost per human genome over time, NHGRI series.
% Data: NHGRI "Cost per Genome", 2001-2023 (snapshot values used here are
% the public series points; full table lives in data/sequencing-costs-nhgri.csv).
\begin{tikzpicture}
  \begin{semilogyaxis}[
      width=12cm, height=8cm,
      xlabel={Year},
      ylabel={Cost per genome (US\$)},
      grid=both,
      xmin=2000, xmax=2024,
      ymin=100, ymax=1e8,
      ymajorgrids=true,
      minor y tick num=9,
      legend style={font=\footnotesize, at={(0.02,0.05)}, anchor=south west},
      title style={font=\small\bfseries},
      title={Cost per human genome, 2001-2023 (NHGRI series)},
    ]
    % NHGRI cost-per-genome series. Values are the project's snapshots
    % (compiled from the public NHGRI graph and tables, see data/README.md).
    \addplot[engblue, thick, mark=*] coordinates {
      (2001, 95000000)
      (2002, 70000000)
      (2003, 50000000)
      (2004, 22000000)
      (2005, 10000000)
      (2006, 7000000)
      (2007, 5000000)
      (2008, 1500000)
      (2009, 350000)
      (2010, 100000)
      (2011, 21000)
      (2012, 8000)
      (2013, 6000)
      (2014, 4000)
      (2015, 1500)
      (2016, 1200)
      (2017, 1100)
      (2018, 1000)
      (2019, 950)
      (2020, 750)
      (2021, 600)
      (2022, 400)
      (2023, 200)
    };
    \addlegendentry{NHGRI cost per genome}
    % Moore's-law reference line, anchored at 2001 = 95M, doubling 24 months
    \addplot[engred, thick, dashed, domain=2001:2023, samples=24]
      {95e6 * 0.5^((x-2001)/2)};
    \addlegendentry{Moore's-law reference ($\times 0.5$ / $2$ y)};
  \end{semilogyaxis}
\end{tikzpicture}

\caption{Cost per human genome from 2001 to 2023 (NHGRI series,
log-scale ordinate). The reference dashed line is Moore's-law
doubling every $24$ months from the 2001 cost of \$95M; sequencing
costs fell faster than Moore's law from $2007$ onward as
next-generation technologies replaced Sanger. The slope flattens
after 2018 as Illumina approached its per-base cost floor; the small
inflection at 2023 reflects the introduction of the Ultima
single-base-extension platform and the Illumina NovaSeq X.}
\label{fig:seq-cost}
\end{figure}

\engterm{Long-read sequencing}. Two technologies (Pacific Biosciences
single-molecule real-time sequencing \cite{paper:v6c03-eid-2009} and
Oxford Nanopore Technologies \cite{paper:v6c03-jain-2018}) produce
reads of $10\,\text{kb}$ to $100\,\text{kb}$ at error rates of
$10^{-2}$ to $10^{-3}$. The long reads resolve repetitive sequences,
structural variants, and full-length transcripts that short-read
sequencing cannot. The combination of long-read scaffolding and
short-read polishing now produces near-complete telomere-to-telomere
human-genome assemblies (the T2T-CHM13 reference of Nurk et al.
\cite{paper:v6c03-nurk-t2t-2022}, 2022, was the first complete human
genome assembled with this approach).

\begin{table}[h]
\centering
\caption{Modern sequencing platforms compared. Throughput, error rate,
and cost per gigabase are order-of-magnitude (current as of 2024);
specific platform models change rapidly and supersede prior generations.}
\label{tab:platforms}
\small
\begin{tabular}{lrrrl}
\toprule
Platform                    & Read length        & Throughput (Gb/run) & Error rate     & Cost (\$/Gb) \\
\midrule
Sanger ABI 3730             & $500$--$1000$ bp   & $0.001$             & $10^{-4}$      & $\sim 500$ \\
Illumina NovaSeq X          & $150$ bp           & $16\,000$           & $10^{-3}$      & $\sim 1.5$ \\
PacBio Revio (HiFi)         & $15$--$25$ kbp     & $90$ per cell       & $10^{-3}$      & $\sim 12$  \\
Oxford Nanopore PromethION  & up to $4$ Mb       & $14\,000$           & $10^{-2}$      & $\sim 5$   \\
\bottomrule
\end{tabular}
\end{table}

\engterm{Engineering applications}. Modern sequencing supports
biomanufacturing (quality control of engineered strains), medicine
(diagnostic genome sequencing, oncology genotyping, infectious-disease
identification), and basic biology (genome-wide association studies,
single-cell transcriptomics, microbiome surveys).
Chapter~\ref{vol06:ch11} on bioinformatics develops the algorithmic
side of sequence analysis; Chapter~\ref{vol06:ch09} develops the
synthetic-biology applications.

\engterm{Cost trajectory and forensic genealogy}. The same cost
curve that supports clinical genomics also supports law-enforcement
applications. The 2018 identification of the Golden State Killer
(Joseph DeAngelo) used a crime-scene DNA profile uploaded to the
public GEDmatch database, identifying distant relatives and triangulating
to a suspect. The technique (forensic investigative genetic genealogy)
relies on the falling cost of sequencing to populate large public
databases of voluntarily-submitted profiles and on genome-wide SNP
data to identify third-cousin matches that point to a small candidate
pool. The case raised the question of whether genealogy-database
users had consented to law-enforcement search of their profiles; the
2019 GEDmatch policy change to opt-in for law-enforcement search and
the 2022 Department of Justice interim guidance on the technique are
the regulatory response in the United States. The engineering reading
is that a public benefit (cold-case resolution) and a private cost
(genomic-privacy erosion for users who did not anticipate the use)
arise from the same falling-cost trajectory; sequencing has reached
the point where the technology is not the constraint and the
constraint shifts to governance.

\engterm{Cell-line authentication and the HeLa cross-contamination
problem}. The 2013 HeLa genome papers \cite{paper:v6c03-landry-2013}
revealed two engineering facts about the line: it is karyotypically
unstable (chromosome counts vary from $70$ to $164$ across substrains)
and it has cross-contaminated dozens of nominally distinct cell lines
in published research. The cross-contamination is detectable by
short-tandem-repeat (STR) profiling, a sequencing-derived technique
that the International Cell Line Authentication Committee now requires
of submissions to major cell-line repositories. The case is a
canonical example of how the same sequencing technology that supports
biomanufacturing also supports the quality-control infrastructure
that catches its own failures.

\section{Failure: replication errors, cancer, hereditary disease}\pathtag{core}

The cell's information-processing machinery fails in characteristic
ways. The failures are mechanism-class, not single named events. We
identify the principal classes.

\engterm{Cancer as a failure of replication control}. Cancer is the
disease of unregulated cell proliferation. Its molecular basis is the
accumulation of mutations in genes that control proliferation, cell
death, and genomic integrity. Two functional classes of cancer gene
are recognised. \engterm{Oncogenes} are normal cell-growth-promoting
genes (proto-oncogenes) whose activated, mutated forms drive
proliferation; activating mutations include point mutations
(\emph{KRAS}, \emph{BRAF}), gene amplifications (\emph{MYC},
\emph{ERBB2}), and chromosomal translocations (the BCR-ABL fusion in
chronic myeloid leukaemia). \engterm{Tumour-suppressor genes} are
normal restraints on proliferation; loss-of-function mutations in
both copies of a tumour suppressor (\emph{TP53}, \emph{RB1},
\emph{APC}, the DNA-mismatch-repair genes) remove the restraint.

\engterm{Knudson's two-hit hypothesis}. Knudson \cite{paper:v6c03-knudson-1971}
analysed the incidence of retinoblastoma, a rare childhood eye tumour
that occurs in two clinical forms: a familial form with bilateral
tumours and a sporadic form with unilateral tumours. He showed that
the familial form follows the kinetics of a single-mutation event
(linear in age, consistent with inheriting one defective \emph{RB1}
allele and acquiring the second somatically) while the sporadic form
follows the kinetics of a two-mutation event (quadratic in age,
consistent with acquiring both \emph{RB1} alleles somatically). The
two-hit hypothesis is the canonical reading of tumour-suppressor
genetics: a tumour suppressor requires loss of both alleles to
release the restraint, so an inherited single-allele loss leaves the
carrier one mutational event short of cancer.

\engterm{Multi-hit cancer progression}. Vogelstein and Kinzler and
their collaborators extended the two-hit reading to a stepwise
multi-hit model of solid-tumour progression
\cite{paper:v6c03-fearon-vogelstein-1990, paper:v6c03-vogelstein-kinzler-2013}.
The canonical colorectal-cancer progression (APC loss in normal
epithelium, KRAS activation in early adenoma, p53 loss in late
adenoma, SMAD4 or DCC loss in carcinoma in situ, additional
chromosomal-instability events in metastasis) is shown in
Figure~\ref{fig:cancer-progression}. A typical solid tumour at
diagnosis carries $10$ to $100$ \engterm{driver} mutations and a
much larger number of \engterm{passenger} mutations that confer no
selective advantage but accumulated by association. Recent large-scale
sequencing efforts (TCGA in the USA, ICGC internationally) have
catalogued the mutation patterns of dozens of cancer types and
identified hundreds of driver genes. The engineering reading is that
cancer is a failure of the cell's proliferation-control system caused
by mutational degradation of the control inputs; the therapeutic
strategies (targeted inhibition of activated oncoproteins, immune
checkpoint blockade to restore immune surveillance, restoration of
damaged tumour-suppressor function) all aim to compensate for the
lost regulatory inputs or to exploit the damaged-control state.

\begin{figure}[h]
\centering
% Knudson two-hit / multi-hit cancer progression: stepwise mutation accumulation.
\begin{tikzpicture}[
    >=Stealth, node distance=2.0cm,
    cell/.style={draw, circle, thick, minimum size=1.2cm, align=center,
      font=\footnotesize, inner sep=1pt},
    normal/.style={cell, fill=englime!30},
    hit/.style={cell, fill=engamber!35},
    tumor/.style={cell, fill=engred!40},
    arrow/.style={->, thick},
    lbl/.style={font=\scriptsize, align=center, fill=white, inner sep=1pt},
  ]
  \node[normal] (s0) {normal \\ APC$^{+/+}$};
  \node[hit, right=of s0] (s1) {APC$^{+/-}$};
  \node[hit, right=of s1] (s2) {APC$^{-/-}$ \\ KRAS$^{+}$};
  \node[hit, right=of s2] (s3) {{+}p53$^{-/-}$};
  \node[tumor, right=of s3] (s4) {SMAD4$^{-}$ \\ +DCC$^{-}$};
  \node[tumor, right=of s4] (s5) {invasive \\ carcinoma};

  \draw[arrow] (s0) -- (s1) node[lbl, midway, above] {APC \\ loss};
  \draw[arrow] (s1) -- (s2) node[lbl, midway, above] {KRAS \\ activate};
  \draw[arrow] (s2) -- (s3) node[lbl, midway, above] {p53 \\ loss};
  \draw[arrow] (s3) -- (s4) node[lbl, midway, above] {SMAD4 \\ DCC};
  \draw[arrow] (s4) -- (s5) node[lbl, midway, above] {invasion};

  % Histology labels under each circle
  \node[font=\footnotesize, below=0.25cm of s0] {epithelium};
  \node[font=\footnotesize, below=0.25cm of s1] {hyperplasia};
  \node[font=\footnotesize, below=0.25cm of s2] {early adenoma};
  \node[font=\footnotesize, below=0.25cm of s3] {late adenoma};
  \node[font=\footnotesize, below=0.25cm of s4] {carcinoma in situ};
  \node[font=\footnotesize, below=0.25cm of s5] {metastasis};

  % Time axis
  \draw[thick, ->] (-0.6, -1.6) -- (14.6, -1.6) node[right] {time (years)};
  \foreach \x/\t in {0/0, 2.9/5, 5.8/10, 8.7/15, 11.6/20, 14.0/30+} {
    \draw (\x, -1.55) -- (\x, -1.65);
    \node[font=\scriptsize, below=2pt] at (\x, -1.65) {\t};
  }

  % Knudson reference
  \node[font=\scriptsize, anchor=west, align=left] at (-0.6, -2.4)
    {Stepwise mutation accumulation: each hit
     confers a selective advantage at the cellular level. \\
     Knudson's two-hit hypothesis (1971) applies to tumour suppressors;
     oncogenes activate by single hits.};
\end{tikzpicture}

\caption{The Fearon-Vogelstein stepwise model of colorectal-tumour
progression. Each step requires both alleles of a tumour suppressor
(APC, p53) or a single activating event on an oncogene (KRAS); the
cumulative chain produces invasive carcinoma over years to decades.
Knudson's two-hit hypothesis applies to each tumour-suppressor step;
the multi-hit composition explains the strongly age-dependent
incidence of solid tumours.}
\label{fig:cancer-progression}
\end{figure}

\engterm{Mismatch-repair failure and the mutator phenotype}. Cells
with defective mismatch-repair machinery (Lynch syndrome
\cite{paper:v6c03-lynch-1985}, microsatellite-unstable colorectal
cancer) have mutation rates $10$ to $1000$ times the normal rate. The
accelerated mutation rate drives faster tumour progression in some
contexts but also produces many neoantigens (mutant proteins
recognised by the immune system), making microsatellite-unstable
tumours unusually responsive to immune checkpoint inhibitors. The
mutator phenotype is an instance of the engineering principle that
loss of a single error-correction stage produces a measurable increase
in the overall error rate, with clinical consequences proportional to
the importance of the lost stage. The categorisation of cancer
genetic instabilities into microsatellite-instability and chromosomal-
instability subtypes is due to Lengauer, Kinzler, and Vogelstein
\cite{paper:v6c03-lengauer-kinzler-vogelstein-1998}; the
mismatch-repair sub-pathway diagnosis (MLH1 promoter methylation
versus germline MMR-gene mutation) determines clinical management.

\engterm{Hereditary disease}. Inherited diseases caused by mutations
in single genes (monogenic disease) follow Mendelian inheritance
patterns: autosomal dominant (\emph{Huntington's} disease,
\emph{Marfan} syndrome), autosomal recessive (cystic fibrosis,
sickle cell disease), X-linked (haemophilia, Duchenne muscular
dystrophy), or mitochondrial (Leber's hereditary optic neuropathy).
The Online Mendelian Inheritance in Man catalogue
\cite{web:v6c03-omim-2024} lists $\sim 6000$ such diseases at the time
of writing. Polygenic diseases (common forms of diabetes, heart
disease, schizophrenia) arise from combinations of many small-effect
variants; whole-genome association studies have identified thousands
of such variants, with each typically contributing a fraction of a
percent to disease risk. The engineering reading is that the cell's
information-processing system can fail at many distinct loci, with
the failure mode at each locus determined by the gene's normal
function. Figure~\ref{fig:punnett} contrasts the inheritance pattern
of an autosomal dominant disease (Huntington's, Hh $\times$ hh) with
an autosomal recessive disease (cystic fibrosis, Cc $\times$ Cc).

\begin{figure}[h]
\centering
% Punnett-square panel: autosomal dominant (Huntington) vs autosomal recessive (CF)
\begin{tikzpicture}[
    >=Stealth,
    cell/.style={draw, thick, minimum width=1.2cm, minimum height=1.2cm,
      align=center, font=\footnotesize, inner sep=1pt},
    affected/.style={cell, fill=engred!35},
    carrier/.style={cell, fill=engamber!30},
    healthy/.style={cell, fill=englime!30},
    title/.style={font=\small\bfseries, align=center},
  ]
  % Left panel: Autosomal dominant, heterozygous x unaffected (Hh x hh)
  \node[title] at (1.8, 4.6) {Autosomal dominant};
  \node[title, font=\footnotesize] at (1.8, 4.15) {Hh $\times$ hh (Huntington)};
  \node[font=\footnotesize] at (0.6, 3.4) {H};
  \node[font=\footnotesize] at (1.8, 3.4) {h};
  \node[font=\footnotesize] at (0.0, 2.8) {h};
  \node[font=\footnotesize] at (0.0, 1.6) {h};
  \node[affected] at (0.6, 2.8) {Hh \\ affected};
  \node[healthy]  at (1.8, 2.8) {hh \\ unaffected};
  \node[affected] at (0.6, 1.6) {Hh \\ affected};
  \node[healthy]  at (1.8, 1.6) {hh \\ unaffected};
  \node[font=\footnotesize] at (1.2, 0.7) {$1{:}1$ affected:unaffected};

  % Right panel: Autosomal recessive, two carriers (Cc x Cc)
  \node[title] at (6.8, 4.6) {Autosomal recessive};
  \node[title, font=\footnotesize] at (6.8, 4.15) {Cc $\times$ Cc (cystic fibrosis)};
  \node[font=\footnotesize] at (5.6, 3.4) {C};
  \node[font=\footnotesize] at (6.8, 3.4) {c};
  \node[font=\footnotesize] at (5.0, 2.8) {C};
  \node[font=\footnotesize] at (5.0, 1.6) {c};
  \node[healthy]  at (5.6, 2.8) {CC \\ unaffected};
  \node[carrier]  at (6.8, 2.8) {Cc \\ carrier};
  \node[carrier]  at (5.6, 1.6) {Cc \\ carrier};
  \node[affected] at (6.8, 1.6) {cc \\ affected};
  \node[font=\footnotesize] at (6.2, 0.7) {$1{:}2{:}1$ healthy:carrier:affected};

  % Legend
  \node[font=\footnotesize, anchor=west] at (-0.4, -0.3) {Legend:};
  \node[healthy] at (1.3, -0.3) {};   \node[font=\scriptsize, anchor=west] at (2.0, -0.3) {unaffected};
  \node[carrier] at (4.3, -0.3) {};   \node[font=\scriptsize, anchor=west] at (5.0, -0.3) {carrier};
  \node[affected] at (6.7, -0.3) {};  \node[font=\scriptsize, anchor=west] at (7.4, -0.3) {affected};
\end{tikzpicture}

\caption{Punnett-square panel. Left: autosomal dominant (one defective
copy is sufficient). One parent carries one disease allele (Hh); the
mating with an unaffected parent (hh) produces $50\%$ affected
offspring. Right: autosomal recessive (two defective copies required).
Two carrier parents (Cc) produce $25\%$ affected offspring and $50\%$
carriers. Population-level incidence follows Hardy-Weinberg from the
carrier frequency in the absence of selection.}
\label{fig:punnett}
\end{figure}

\engterm{Three canonical hereditary diseases}. Sickle cell disease
(autosomal recessive, HBB $\beta^6$ Glu$\to$Val substitution) provides
the classic example of a single-codon mutation producing a major
phenotype: the substitution causes haemoglobin polymerisation under
hypoxia, deforming red blood cells. The disease frequency in
sub-Saharan Africa ($\sim 1.5\%$ of births in some populations) is
maintained by carrier resistance to \emph{Plasmodium falciparum}
malaria, a balancing-selection example
\cite{web:v6c03-genereviews-2024}. Cystic fibrosis (autosomal
recessive, CFTR; most commonly the $\Delta$F508 in-frame deletion)
causes loss of chloride-channel function in epithelial cells, with
clinical consequences in lung, pancreas, and gut; the $\sim 3.5\%$
carrier frequency in northern European populations is one of the
highest known carrier frequencies for a serious autosomal recessive
disease. Huntington's disease (autosomal dominant, HTT CAG-repeat
expansion) provides the trinucleotide-repeat-expansion archetype: an
unstable CAG repeat in exon 1 of HTT expands across generations,
with disease appearing when the repeat exceeds $\sim 40$ copies and
age of onset inversely correlated with repeat length. The repeat-
expansion mechanism breaks the standard expectation that mutations
affect only the carrier; expansions can occur during gametogenesis,
particularly in spermatogenesis, so paternal inheritance is associated
with earlier-onset offspring.

\engterm{Replication errors during ageing}. The cells of older
animals carry more mutations than the cells of younger animals.
Mutations accumulate at a rate that integrates the per-division
error rate, the cell-division frequency, and the contribution from
unrepaired DNA damage. The accumulation is implicated in many of the
chronic diseases of ageing (cancer, neurodegeneration,
cardiovascular disease) and is one of the volume's running themes
on cellular failure \cite{paper:v6c03-blokzijl-2016}.
Chapter~\ref{vol06:ch07} returns to ageing in the context of immune
senescence; the present chapter limits itself to the genetic-
information side.

\engterm{Named case: HeLa cells}. HeLa cells, derived in 1951 from
a cervical-tumour biopsy of Henrietta Lacks at Johns Hopkins
Hospital, became the first continuously propagated human cell line
and have served as the workhorse of cell biology for seven decades.
The biopsy was taken without consent, in accordance with the standard
practice of the era, and the Lacks family learned of the cell line
only in 1973. The case is the canonical example of a failure of
informed consent in research that produced both major public benefit
(polio-vaccine development, the discovery of the HPV-cervical cancer
link, the recognition of telomerase as the immortalisation enzyme)
and a long-latency ethical reckoning that culminated in the 2013
NIH agreement with the Lacks family creating controlled-access
review of HeLa genome data \cite{paper:v6c03-skloot-2010,
paper:v6c03-landry-2013, paper:v6c03-hudson-collins-2013,
acc:hela-lacks-1951}. The mechanism is engineering-relevant in two
ways. First, the line is karyotypically unstable (chromosome counts
from $70$ to $164$ across substrains), so the cells in any given
laboratory are not the cells of any other; the implication for
reproducibility is documented in the cell-line-authentication
literature. Second, publishing a cell line's whole-genome sequence
discloses heritable information about the donor's living relatives, a
genomic-privacy problem that has no analogue in pre-genomic biology;
the 2013 NIH-Lacks agreement is the first attempt at a workable
framework for the disclosure question.

\engterm{Named case: He Jiankui CRISPR babies}. He Jiankui announced
in November 2018 the birth of twin girls from human embryos he had
edited with CRISPR-Cas9 to disrupt the CCR5 gene for putative HIV
resistance \cite{paper:v6c03-cyranoski-ledford-2018,
paper:v6c03-musunuru-2019, paper:v6c03-greely-2019,
paper:v6c03-who-genome-editing-2021, acc:he-jiankui-crispr-2018}. The
case is a multi-layer failure. The technical mechanism failed: the
edits did not reproduce the natural CCR5-$\Delta$32 deletion that
confers HIV resistance, the embryos were mosaic (different cells in
the same embryo carrying different edits), and off-target risk at
homologous loci was not adequately characterised before implantation.
The clinical rationale failed: sperm-washing protocols reduce vertical
HIV transmission risk to near zero in the clinical setting that
the parents faced, so no plausible benefit balanced the unknown
off-target risk. The regulatory mechanism failed: He bypassed the
Chinese Ministry of Health guidance on heritable human genome
editing; ethics-committee documentation was partly forged; parental
consent did not adequately describe the experimental nature of the
work. He was convicted in December 2019 of illegal medical practice
and sentenced to three years' imprisonment. The case has become the
canonical example of a normative-consensus-without-enforcement
failure in biotechnology governance and has driven the WHO 2021
governance framework and the International Commission's 2020 report
establishing the requirement that heritable human genome editing not
proceed until standards for safety, efficacy, and societal
authorisation are in place. The engineering reading is that
laboratory capability ran ahead of governance capability and that
the absence of binding international standards (the prior norms
were advisory only) was the dominant single failure mode.

\begin{archetype}[Information, failure, and scaling]
The cell's information system exhibits three of the volume's
archetypes in their canonical biological forms. The
\engterm{information archetype} is the digital storage and
controlled retrieval of sequence information at the molecular level,
with error rates corresponding to those of engineered storage. The
\engterm{failure archetype} is the cascade of failure modes
(replication error, repair failure, regulatory loss, proliferation
loss) that produce cancer, hereditary disease, and ageing-associated
pathology. The \engterm{scaling archetype} is the way the same
machinery, sized differently, supports genomes from $4 \times 10^3$
base pairs (small viroids) to $10^{11}$ base pairs (large plant
genomes) with similar per-base error rates and similar architectural
features.
\end{archetype}

\begin{estimation}
\textbf{Question.} Estimate the time it takes to read an entire human
genome on a modern short-read sequencer.

\smallskip
\textit{Estimate, before reading on.}

\smallskip
The human haploid genome is $3.2 \times 10^9\,\text{bp}$. To obtain
reliable variant calls, $30\times$ coverage is the working standard:
each base is sequenced an average of $30$ times in independent reads.
The total sequenced bases are $9.6 \times 10^{10}\,\text{bp}$.

A modern Illumina NovaSeq X (current as of 2024) produces approximately
$1.6 \times 10^{13}\,\text{bp}$ of output per $48\,\text{hour}$ run
\cite{paper:v6c03-illumina-novaseq-2024}. The rate is
$1.6 \times 10^{13} / (48 \times 3600) = 9.3 \times 10^7\,
\text{bp}/\text{s}$.

The time to sequence one $30\times$ human genome:
\[
T = \frac{9.6 \times 10^{10}}{9.3 \times 10^7} \approx 1030\,\text{s} \approx 17\,\text{minutes}.
\]

In practice the instrument multiplexes many samples per run (the
$48\,\text{h}$ run produces approximately $160$ human genomes at
$30\times$ each), so the per-sample wall-clock time is a fraction of
the run time but the throughput per sample is the figure above.

\smallskip
\textit{Postmortem.} The estimate is the throughput limit, not the
end-to-end time. Sample preparation, library construction, quality
control, base calling, alignment, and variant calling each add
time. The actual turn-around for a clinical genome sequencing is
typically $1$ to $2\,\text{weeks}$ at the time of writing, dominated
by sample handling rather than by the instrument itself. The
sequencing throughput has reached the point where it is no longer
the rate-limiting step; the rate-limiting step is now the
informatics, the interpretation, and the clinical reporting. The
shift identifies the next engineering bottleneck and the focus of
Chapter~\ref{vol06:ch11}.
\end{estimation}

\begin{estimation}
\textbf{Question.} Estimate the Illumina sequencing throughput needed
to call rare somatic variants at $0.1\%$ allele frequency with
$30\times$ mean coverage, and compare to the $30\times$ germline
standard.

\smallskip
\textit{Estimate, before reading on.}

\smallskip
To call a single-nucleotide variant present in $0.1\%$ of reads at a
given position, the standard rule of thumb is to require at least
$5$ supporting reads carrying the variant. The expected number of
variant reads at coverage $C$ is $C \cdot 0.001$. Setting $C \cdot
0.001 = 5$ gives $C = 5000$, an order-of-magnitude higher than the
$30 \times$ germline standard. The actual workable depth is higher
once one accounts for sequencing error (Illumina substitution error
rate $\sim 10^{-3}$ per base; a depth of $5000$ would produce $5$
false positives per position from sequencing error alone, exactly
matching the true-positive signal), so a target depth in the
$10\,000$ to $30\,000$ range is typical for liquid-biopsy cell-free
DNA assays, often combined with unique molecular identifiers (UMIs)
to suppress sequencing-error false positives.

The total sequenced bases per human genome at $30\,000 \times$ coverage
are $9.6 \times 10^{13}$, or $1000 \times$ the germline-genome cost.
A $\$200$ germline genome implies a $\sim \$200\,000$ deep
liquid-biopsy assay at full genome scale; in practice the
liquid-biopsy industry targets a small panel of $\sim 100$ to
$500$ cancer-relevant genes ($\sim 1\,\text{Mb}$ of target) at
$30\,000 \times$ depth, bringing the per-sample cost to
$\sim \$1000$ to $\$3000$. The estimation identifies the
single dimension (variant allele frequency) along which sequencing
depth requirements scale and the engineering response
(targeted panels plus UMIs) that brings the costs back into clinical
range.

\smallskip
\textit{Postmortem.} The same estimation applied to viral or
microbial variants at $0.01\%$ frequency demands depths in the
$10^5$ range, putting whole-genome sequencing out of reach for
routine application and identifying targeted-amplicon sequencing as
the practical method for low-frequency-variant detection. The depth
budget at any given allele frequency is the single most useful
back-of-envelope for sequencing experimental design.
\end{estimation}

\section{Project}\pathtag{core}

\begin{project}[Simulate genetic drift in a small population]
The reader simulates the Wright-Fisher model of genetic drift in a
finite population, observes the time-to-fixation distribution, and
compares to the analytic predictions.

\textbf{Track.} Simulation.

\textbf{Hazard class}: none.

\textbf{Model.} A diploid population of $N$ individuals, with a single
locus carrying two alleles ($A$ and $a$). Each generation, $2N$
gametes are sampled with replacement from the previous generation's
allele pool. Initial allele frequency $p_0 = 0.5$.

\textbf{Procedure.}

\begin{enumerate}[leftmargin=*]
\item Implement the Wright-Fisher sampling step in a programming
  language of your choice. A starter implementation is provided at
  \texttt{code/wright\_fisher\_drift.py}.
\item For population sizes $N \in \{10, 100, 1000\}$, simulate
  $1000$ independent populations and record the generation at which
  each population fixes one allele or the other.
\item Plot the empirical distribution of fixation times and compute
  the mean.
\item Compare with the analytic prediction that the expected time to
  fixation conditional on $p_0 = 0.5$ scales as $\sim 2.8 N$ generations.
\item Repeat with a selective advantage $s = 0.01$ for the $A$ allele
  and report how the probability of $A$ fixation depends on $N s$. The
  Kimura diffusion result $P_\text{fix}(p_0) =
  (1 - e^{-2 N s p_0})/(1 - e^{-2 N s})$ is the analytic reference.
\item \emph{Extension}. Compute the effective population size
  $N_e$ for a sampled real-population SNP frequency trajectory by
  fitting the variance accumulation $\text{Var}[p(t)] \approx
  p_0(1-p_0)(1 - e^{-t/(2 N_e)})$; comment on how $N_e$ is generally
  smaller than the census population size.
\end{enumerate}

\textbf{Deliverable.} A five- to ten-page report with one
fixation-time-distribution figure, one fixation-probability-versus-$N s$
figure, source code, and a paragraph identifying which analytic
prediction your simulation reproduces most accurately and which least.

\textbf{Time.} Four to eight hours.
\end{project}

\section{Exercises}

\subsection*{Calculation}\pathtag{core}

\begin{exercise}
The human haploid genome is $3.2 \times 10^9\,\text{bp}$. Compute
the information content in bits and in bytes. Compare to the storage
capacity of a typical USB flash drive (current as of 2024).
\end{exercise}

\paragraph{Solution sketch.} At $\log_2 4 = 2$ bits per base pair,
the raw content of the haploid genome is
$2 \cdot 3.2 \times 10^9 = 6.4 \times 10^9$ bits, or
$6.4 \times 10^9 / 8 = 8.0 \times 10^8$ bytes $= 800\,\text{MB}$.
The diploid value is $1.6\,\text{GB}$. A typical 2024-era USB flash
drive holds $64$ to $256\,\text{GB}$, so $80$ to $320$ haploid
genomes uncompressed. Lossless compression of human DNA (removing
repetitive sequence such as Alu and L1 elements) reduces the
content to $\sim 400$ to $700\,\text{MB}$ per haploid genome; the
compressed representation used by the SAM/BAM/CRAM file format
brings a $30\times$-coverage sequenced genome (typical $90\,\text{GB}$
raw) down to $30$ to $40\,\text{GB}$ on disk, so a $128\,\text{GB}$
USB drive can carry approximately three sequenced human genomes
with read data, or many tens of compressed reference sequences.

\begin{exercise}
DNA polymerase has a base-incorporation rate of approximately
$1000\,\text{bp}/\text{s}$ in \emph{E.~coli}. Compute the time
required to replicate the entire \emph{E.~coli} genome
($4.6 \times 10^6\,\text{bp}$) from a single bidirectional origin.
Compare to the observed doubling time of $30\,\text{min}$.
\end{exercise}

\paragraph{Solution sketch.} A bidirectional fork from a single origin
moves two forks at $1000$ bp/s each, so each fork must traverse half
the genome: $2.3 \times 10^6$ bp at $1000$ bp/s
$= 2300\,\text{s} \approx 38\,\text{minutes}$. The observed doubling
time is $30\,\text{min}$. The replication time (alone) is therefore
\emph{larger} than the cell-division cycle. The resolution: \emph{E.
coli} initiates a new round of replication before the previous round
completes (overlapping rounds), so the chromosome carries multiple
replication forks simultaneously during fast growth. The next-generation
chromosome inherits a partially replicated chromosome and continues.
The phenomenon is called ``multifork replication'' and is the
engineering trick that makes $30$-minute doubling consistent with a
single-origin $40$-minute replication time. The number of overlapping
rounds at steady state is roughly $\tau_C / \tau_D$ where $\tau_C$ is
the chromosome-replication time and $\tau_D$ is the doubling time.

\begin{exercise}
The persistence length of DNA is $L_p = 50\,\text{nm}$. Compute the
number of base pairs in one persistence length and the bending energy
required to curve a DNA segment of length $L_p$ into a complete
circle ($U = 2\pi^2 k_B T L_p/L$). Use $T = 310\,\si{\kelvin}$.
\end{exercise}

\paragraph{Solution sketch.} At $0.34\,\text{nm}$ per base pair,
$50\,\text{nm}$ is $50/0.34 \approx 147$ bp, almost exactly the
nucleosomal wrap of $147$ bp around the histone core. The energy to
bend a segment of length $L$ into a circle of circumference $L$ is
$U = 2\pi^2 k_B T L_p / L$. Setting $L = L_p = 50\,\text{nm}$:
$U = 2\pi^2 \cdot k_B T \cdot 1 = 19.7\,k_B T$. With
$k_B T = 4.28 \times 10^{-21}\,\text{J}$ at $T = 310\,\text{K}$,
$U = 8.4 \times 10^{-20}\,\text{J} \approx 0.52\,\text{eV}$. The
energy is large compared to thermal fluctuations ($k_B T$) and
substantial compared to a hydrogen bond ($\sim 5$--$20\,k_B T$); a
single histone-DNA interaction of comparable magnitude is what
stabilises the wrap. The matched numbers (one persistence length
equals one nucleosome equals $\sim 20\,k_B T$ of bending energy) are
not coincidence: the cell's chromatin architecture is built to
exactly match the polymer's mechanical scale.

\begin{exercise}
Compute the energetic cost (in ATP equivalents) of synthesising a
mRNA of $1000\,\text{nucleotides}$ given the cost per nucleotide of
$2$ ATP equivalents (one for the triphosphate bond, one for the
nucleotide diphosphate synthesis). Compute the cost of translating
the $333$-amino-acid protein it encodes.
\end{exercise}

\paragraph{Solution sketch.} Transcription: $1000\,\text{nt} \times
2\,\text{ATP/nt} = 2000$ ATP equivalents. Translation: each peptide
bond consumes $4$ ATP equivalents ($2$ for the aminoacyl-tRNA
charging, $1$ for the elongation factor EF-Tu cycle, $1$ for
translocation), so a $333$-amino-acid protein costs
$333 \times 4 = 1332$ ATP for elongation, plus a one-time initiation
cost of $\sim 2$ ATP and a termination cost of $\sim 1$ ATP. Total
$\approx 1335$ ATP for translation, $\sim 3300$ ATP for the full
mRNA-to-protein cost of one protein copy. With $30\,\text{kJ/mol}$
per ATP hydrolysis, the per-protein synthesis cost is
$\sim 10^5\,\text{kJ/mol}$ of protein, or roughly $1.6 \times
10^{-19}\,\text{J}$ per molecule. A bacterium at $\sim 10^6$ protein
copies per generation expends $\sim 10^{-13}\,\text{J}$ on
protein synthesis, comparable to the total metabolic budget per
division.

\begin{exercise}
The mutation rate in \emph{E.~coli} is approximately $10^{-9}$ per
base pair per generation. Compute the expected number of mutations
per cell per generation in the $4.6 \times 10^6\,\text{bp}$ genome.
Compute the expected number of mutations per gene per generation for
a typical $1000\,\text{bp}$ gene.
\end{exercise}

\paragraph{Solution sketch.} Per-cell per-generation:
$10^{-9} \times 4.6 \times 10^6 = 4.6 \times 10^{-3}$ mutations, or
one new mutation every $\sim 220$ generations on average. Per typical
$1000\,\text{bp}$ gene: $10^{-9} \times 10^3 = 10^{-6}$ per generation,
or one mutation every $10^6$ generations of that single gene. With
a population size of $10^9$ cells in an overnight \emph{E. coli}
culture, the gene experiences approximately $10^3$ independent
mutational events per overnight (one every $10^6$ generations $\times
10^9$ independent cell lines), enough to seed every possible
single-nucleotide variant of that gene many times over in a single
laboratory experiment. The implication is that mutational saturation
of a small target is routine in bacterial selection experiments and
mutational coverage is rarely the limiting factor.

\begin{exercise}
The human germline mutation rate is approximately
$1.2 \times 10^{-8}$ per base pair per generation. Compute the
expected number of new mutations in a child relative to the parents
in the $6.4 \times 10^9\,\text{bp}$ diploid genome.
\end{exercise}

\paragraph{Solution sketch.} $1.2 \times 10^{-8} \times 6.4 \times
10^9 = 77$ new mutations per child. The figure agrees with the Kong
et al. trio-sequencing result of $\sim 70$ new variants per child
\cite{paper:v6c03-kong-2012}. Of these, approximately $1$ falls in a
protein-coding region (since coding sequence is $\sim 1.5\%$ of the
genome) and roughly $0.5$ to $1$ is a non-synonymous coding change.
Fewer than $1$ in $1000$ such non-synonymous changes is expected to
be strongly pathogenic; the typical child carries no de novo
strongly-pathogenic coding variant. The arithmetic is the basis for
the engineering claim that human germline mutation, at its measured
rate, would take $\sim 10^4$ generations to accumulate enough
deleterious load to compromise the average protein-coding gene,
matching the timescale at which observable speciation occurs.

\begin{exercise}
A polymerase has an intrinsic error rate of $10^{-5}$ from base
selection and proofreading reduces it to $10^{-7}$. Mismatch repair
catches $99\,\%$ of the remaining errors. Compute the effective
error rate at fixation.
\end{exercise}

\paragraph{Solution sketch.} The composition of the three stages is
the product of their residual rates. After base selection: $10^{-5}$.
After proofreading: $10^{-7}$ (a $100$-fold reduction, residual factor
$0.01$). After mismatch repair: $10^{-7} \cdot (1 - 0.99) = 10^{-9}$.
The fixation error rate is $10^{-9}$ per base pair, matching the
measured rate for fast-growing bacteria. The point of the exercise is
the multiplicative composition: each stage reduces the residual rate
by a factor set by its molecular fidelity, and the overall rate is
the product. Loss of any single stage produces a clinically detectable
mutator phenotype (proofreading loss $\to$ POLE-mutated cancers;
MMR loss $\to$ Lynch syndrome).

\subsection*{Derivation}\pathtag{standard}

\begin{exercise}
Derive the expected time to fixation by genetic drift in a haploid
Wright-Fisher population of size $N$, starting from allele frequency
$p_0$, conditional on eventual fixation of the focal allele.
\end{exercise}

\paragraph{Solution sketch.} The Wright-Fisher process is a Markov
chain on the allele-count state space $\{0, 1, \ldots, N\}$ with
$0$ and $N$ absorbing. In the diffusion limit (large $N$, small
$1/N$) the allele frequency $p(t)$ satisfies
$dp = \sqrt{p(1-p)/N}\,dW$. The unconditional expected fixation time
from $p_0$ is, by Kimura's diffusion result,
$T(p_0) = -4 N [p_0 \ln p_0 + (1 - p_0) \ln(1 - p_0)]$ (for the
diploid generalisation; halve for haploid). At $p_0 = 0.5$ this gives
$T(0.5) = -4 N [0.5 \ln 0.5 + 0.5 \ln 0.5] = 4 N \ln 2 \approx
2.77 N$. Conditional on fixation of the focal allele (i.e.\ throwing
out the runs that lose it), the conditional expected fixation time
is
$T_\text{fix}(p_0) = -[2/(1 - p_0)] N \cdot (1 - p_0) \ln(1 - p_0)
= -2 N \ln(1 - p_0)$ in the limit of small $p_0$; at $p_0 = 0.5$
it equals the unconditional $\sim 2.8 N$. The factor $2.8$ is the
benchmark against which the chapter's simulation script reports
its empirical mean.

\begin{exercise}
Derive the Hardy-Weinberg equilibrium for diploid allele frequencies
under random mating and no selection, mutation, or migration.
\end{exercise}

\paragraph{Solution sketch.} Let allele $A$ have frequency $p$ and
allele $a$ frequency $q = 1 - p$. Under random mating, the offspring
genotype frequencies are the product of allele frequencies in the
parental gametes: $P(AA) = p^2$, $P(Aa) = 2pq$, $P(aa) = q^2$. The
allele frequency in the offspring generation is
$p' = P(AA) + \frac{1}{2}P(Aa) = p^2 + pq = p(p + q) = p$. The allele
frequency is unchanged, and one generation of random mating is
sufficient to bring genotype frequencies into the
$\{p^2, 2pq, q^2\}$ equilibrium. The equilibrium is the null model
against which departures (selection, assortative mating, population
structure, inbreeding) are tested. Practical use: the carrier
frequency for an autosomal recessive disease with prevalence
$10^{-4}$ ($q^2 = 10^{-4}$) is $2 p q \approx 2 q = 0.02$, i.e.\ $1$
in $50$, far higher than the prevalence and a key counselling input
for population carrier screening.

\begin{exercise}
Derive the expected number of synonymous and non-synonymous
substitutions per codon under a uniform random-codon-substitution
model, given the genetic-code table. Use the result to motivate the
$d_N/d_S$ statistic used in molecular-evolution studies.
\end{exercise}

\paragraph{Solution sketch.} Consider a codon at a fixed amino-acid
position. Each codon has nine possible single-nucleotide neighbours
(three positions $\times$ three alternative bases). For each codon
$c$, count how many of the nine neighbours encode the same amino acid
($n_S(c)$, synonymous neighbours) and how many encode a different
amino acid or a stop codon ($n_N(c) = 9 - n_S(c)$, non-synonymous +
nonsense). Averaging over the codon frequencies in a reference set,
the expected fraction of single-nucleotide substitutions that are
synonymous is $\bar{n}_S / 9$ and the expected fraction non-synonymous
is $\bar{n}_N / 9$. For a uniform codon distribution over the standard
code, the average is $\bar{n}_S \approx 2.2$ and $\bar{n}_N \approx
6.8$, so synonymous changes are $24\%$ and non-synonymous $76\%$ of
all single-base mutations under neutral drift. The $d_N/d_S$ statistic
normalises the observed number of non-synonymous substitutions per
non-synonymous \emph{site} ($d_N$) by the analogous synonymous quantity
($d_S$); under neutral drift the ratio is $1$, under purifying
selection $d_N/d_S < 1$, under positive selection $d_N/d_S > 1$. The
codon-table arithmetic is the calibration that makes the ratio
interpretable.

\subsection*{Estimation}\pathtag{core}

\begin{exercise}
Estimate the total mass of DNA in the human body, given that there
are approximately $4 \times 10^{13}$ cells each containing
$\sim 6\,\text{pg}$ of DNA.
\end{exercise}

\paragraph{Solution sketch.} Total DNA mass:
$4 \times 10^{13}\,\text{cells} \times 6 \times 10^{-12}\,\text{g/cell}
= 240\,\text{g}$. A reasonable adult-human DNA mass is roughly a
quarter of a kilogram, distributed across all somatic cells. The
estimation is the input to several follow-on calculations: the total
DNA-base count is $\sim 2 \times 10^{23}$, the molar weight gives
$\sim 0.7\,\text{mol}$ of nucleotides, and the total information
content if every cell's genome were independent would be $\sim 200$
exabytes (i.e.\ $200 \times 10^{18}$ bytes). The figure usefully
contrasts with the $0.8\,\text{GB}$ per cell: the total body holds
a copy of the genome in every cell, an extreme redundancy in
information terms but a structural necessity in biological terms (the
chromatin is the substrate the cell reads, not a database).

\begin{exercise}
Estimate the rate at which a typical cell repairs oxidative DNA
damage, given that approximately $10\,000$ oxidative lesions occur
per cell per day and a healthy cell maintains a steady-state count
of fewer than $100$ unrepaired lesions.
\end{exercise}

\paragraph{Solution sketch.} At steady state, the lesion-formation
rate $k_\text{form}$ equals the lesion-repair rate $k_\text{rep}
\cdot N_\text{lesions}$, so $k_\text{rep} = k_\text{form} /
N_\text{lesions}$. With $k_\text{form} = 10\,000$ lesions/day and
$N_\text{lesions} = 100$, the per-lesion repair rate is
$100\,\text{lesion}^{-1}\,\text{day}^{-1}$, i.e.\ each lesion is
repaired with a characteristic time of $\sim 1/100\,\text{day} =
\sim 14\,\text{minutes}$. The order-of-magnitude agreement with
direct measurements of 8-oxoG removal half-life (minutes to hours)
is the cross-check. The cell repairs oxidative damage with a turnover
time short compared to the cell cycle, which is why a single cell
division contributes negligibly to lifetime mutational load \emph{at
homeostasis}; ageing and disease shift the kinetics by loading the
formation rate, slowing the repair rate, or both, and the steady-state
lesion count rises accordingly.

\begin{exercise}
Estimate the number of distinct human protein-coding gene sequences
that are perfectly preserved across all $8 \times 10^9$ people alive
today (current as of 2024). Use the per-base mutation rate and the
expected accumulated mutation count per gene.
\end{exercise}

\paragraph{Solution sketch.} The average human protein-coding gene
is approximately $1500\,\text{bp}$ of coding sequence. Across $8
\times 10^9$ people, each gene exists in $\sim 1.6 \times 10^{10}$
copies (diploid). The probability that a given $1500$-bp gene copy
is mutation-free at any single base is $1 - \mu \approx 1$ for
$\mu = 10^{-8}$, but across $1500$ bases the probability of being
free of new mutations (since the most recent common ancestor) is
$(1 - \mu)^{1500} \approx e^{-1.5 \times 10^{-5}} \approx 1 - 1.5
\times 10^{-5}$ per generation. Accumulating over many generations
since the MRCA (taking $\sim 10\,000$ generations as a rough scale
for human-population MRCA depth), the per-generation rate compounds
to $\sim 0.15$ accumulated coding mutations per gene per generation,
so practically \emph{no} gene is perfectly preserved in sequence
across the entire population at the time of writing. The standing
population variation is the natural state. The estimate reframes
``preservation'' from sequence identity to functional identity: the
fraction of genes whose every variant in the standing population is
functionally equivalent is what genetic-population surveys actually
measure.

\begin{exercise}
Estimate the cost (in dollars per base pair) of writing a designed
DNA sequence by current commercial gene-synthesis services (current
as of 2024). Compare to the cost per base pair of reading by
short-read sequencing.
\end{exercise}

\paragraph{Solution sketch.} Commercial double-stranded gene synthesis
in 2024 prices at roughly $\$0.05$ to $\$0.10$ per base pair for
standard $0.5$ to $5\,\text{kb}$ fragments (current as of 2024).
Read-by-sequencing on Illumina platforms is now under $\$0.003$ per
megabase, or $\$3 \times 10^{-9}$ per base pair (Table~\ref{tab:platforms}).
The write-to-read cost ratio is therefore $\sim 10^7$: writing a base
costs ten million times what reading one costs. The gap is the
single most informative input for understanding why ``DNA storage''
demonstrations price gigabyte-scale archival storage in millions of
dollars and read it back for tens of dollars. The gap is closing at
$\sim 30\%$ per year on the write side (enzymatic synthesis platforms
expected to displace phosphoramidite chemistry over the next decade),
but the read-write asymmetry has been the dominant feature of the
field for two decades.

\subsection*{Simulation}\pathtag{mastery}

\begin{exercise}
Simulate the accumulation of point mutations in a single
$3 \times 10^9\,\text{bp}$ genome over $10^4$ generations at the
mammalian mutation rate. Plot the mutation count per gene over time
and identify the variance.
\end{exercise}

\paragraph{Solution sketch.} With $\mu = 10^{-8}$ per bp per
generation and a $1500$-bp coding gene, the per-gene per-generation
mutation rate is $1.5 \times 10^{-5}$. Over $10^4$ generations the
expected number of mutations per gene is $0.15$, Poisson-distributed
with variance $0.15$. A simulation drawing from this Poisson at each
generation and accumulating across all $20\,000$ protein-coding genes
gives an expected $0.15 \times 20\,000 = 3000$ mutations per
$10^4$-generation lineage. The variance in mutation count across
genes (varying gene length and base composition) is dominated by
length variation: long genes (e.g.\ titin at $\sim 100\,\text{kb}$
coding sequence) accumulate $\sim 100 \times$ more mutations than
average. The simulation reproduces the empirical observation that
mutational accumulation across the human exome is roughly proportional
to gene length, with corrections for chromatin state and replication
timing of order $30\%$ at most.

\begin{exercise}
Implement a simulation of the Luria-Delbr\"uck fluctuation test for a
bacterial population. Show that the variance of mutant counts across
parallel cultures distinguishes random spontaneous mutation
(Luria-Delbr\"uck distribution) from induced mutation (Poisson
distribution).
\end{exercise}

\paragraph{Solution sketch.} Simulate $K$ parallel cultures, each
starting from $N_0$ wild-type cells and dividing $g$ generations to
reach a final population of $N_0 2^g$ cells. At each division, each
cell mutates with probability $\mu = 10^{-9}$. Track the number of
mutants $X_k$ in culture $k$ at the final generation. Under random
mutation, the mutant count follows the Luria-Delbr\"uck distribution
\cite{paper:v6c03-luria-delbruck-1943}: mutations occurring early in
the culture's history produce large mutant clones (jackpots), so
the distribution has very high variance, with $\text{Var}[X] /
\E[X] \gg 1$. Under the alternative hypothesis (mutations induced
by the selective agent), each cell mutates independently at the time
of selection, so the mutant count is Poisson with $\text{Var}[X] /
\E[X] = 1$. The observed variance-to-mean ratio
$\gg 1$ in Luria and Delbr\"uck's 1943 experiment was the decisive
evidence for spontaneous, pre-selection mutation. The simulation
reproduces both distributions and is the canonical illustration of
how a statistical signature distinguishes two mechanistic
hypotheses.

\subsection*{Design}\pathtag{standard}

\begin{exercise}
Design an error-correcting DNA storage system that encodes
$1\,\text{MB}$ of arbitrary digital information into a synthetic DNA
sequence. Specify the encoding scheme, the redundancy level, and the
expected retention time under refrigerated storage. Compare to
current commercial DNA-storage demonstrations.
\end{exercise}

\paragraph{Solution sketch.} A $1\,\text{MB}$ payload is
$8 \times 10^6$ bits. At $2$ bits per nucleotide, the raw payload is
$4 \times 10^6$ nucleotides. Practical encodings avoid long
homopolymer runs (which are error-prone in synthesis and reading)
and balance GC content within a target window of $40$--$60\%$. A
sensible design: fragment the payload into $200$-nt oligos
($20\,000$ fragments), prepend each with a $20$-nt address index
plus a $20$-nt random barcode for deduplication, and add Reed-Solomon
outer encoding at $20\%$ redundancy plus inner Hamming codes for
intra-oligo correction. Total synthesised length is approximately
$5 \times 10^6$ nucleotides. At $\$0.05$/nt the synthesis cost is
$\$250\,\text{k}$ per $1\,\text{MB}$ archive. Readback by short-read
sequencing at $1000 \times$ coverage takes $\sim 1$ minute on a
NovaSeq X. Retention at $\degC{20}$ refrigeration is decades; in
encapsulated form at room temperature with desiccation it extends
to centuries. Commercial demonstrations (Microsoft / University of
Washington, Twist Bioscience / Catalog Technologies, current as of
2024) operate at the few-MB scale and report similar designs. The
exercise is calibration of the write-cost barrier and the
error-correction overhead.

\begin{exercise}
Design a sequencing strategy for the genome of a previously
unsequenced eukaryote of $10^9\,\text{bp}$. Specify the mix of
short-read and long-read sequencing, the depth, and the expected
cost (current as of 2024 prices). Identify which steps of the
assembly are likely to be hardest.
\end{exercise}

\paragraph{Solution sketch.} A standard 2024 hybrid-assembly recipe
for a $1\,\text{Gb}$ eukaryotic genome: $30\times$ PacBio HiFi reads
($\sim 30\,\text{Gb}$ at $\$12$/Gb $= \$360$); $60\times$ Illumina
short reads ($\sim 60\,\text{Gb}$ at $\$1.50$/Gb $= \$90$);
$30\times$ Hi-C contact data for scaffolding ($\sim \$1500$); and
optical-mapping data for $\sim \$2000$ if structural-variant
resolution is required. Total sequencing cost approximately
$\$4000$ to $\$5000$. The hardest assembly step is typically
resolving large repetitive arrays (centromeric satellites,
heterochromatic blocks); the second-hardest is correctly phasing
heterozygous haplotypes in highly polymorphic regions (the major
histocompatibility complex in vertebrates). Annotation (predicting
gene locations, splicing structure, regulatory elements) routinely
costs more in person-hours than sequencing in dollars: a competent
annotation pipeline for a novel genome takes several months of
bioinformatic labour even on assembled sequence. The estimate
identifies where the bottleneck has shifted: sequencing is now the
cheapest line item; the assembly and annotation labour is the rate-
limiting cost.

\begin{exercise}
Design a CRISPR-Cas9 guide RNA for editing a specific position in the
human genome. Specify the constraints on guide RNA design (length,
target specificity, PAM sequence) and the off-target risks to be
checked.
\end{exercise}

\paragraph{Solution sketch.} For \emph{Streptococcus pyogenes} Cas9
\cite{paper:v6c03-doudna-charpentier-2014}, design a $20$-nt guide
RNA matching the genomic target strand, with the $3$-nt
protospacer-adjacent motif (PAM) sequence \texttt{NGG}
immediately downstream of the target on the non-targeted strand.
Quality constraints: avoid guides with very high or low GC content
(target $40$--$60\%$); avoid runs of four or more identical
nucleotides; avoid guides whose seed sequence (the PAM-proximal
$10$--$12$ nt) matches multiple genome locations. Off-target risk is
the primary safety concern: search the candidate guide against the
reference genome allowing up to $3$ to $4$ mismatches, weight by
mismatch position (PAM-proximal mismatches more disruptive than
PAM-distal), and exclude guides with high-risk off-target hits in
coding regions or essential genes. Verification: amplicon sequence
the on-target site and the top $5$ predicted off-target sites in
edited cells to quantify the edit efficiency and the off-target rate.
The He Jiankui case is the canonical illustration of what happens
when this verification is not performed before clinical application
\cite{acc:he-jiankui-crispr-2018}.

\subsection*{Diagnosis and reverse engineering}\pathtag{core}

\begin{exercise}
A patient's tumour exhibits microsatellite instability and an
unusually high mutation rate. Identify the likely failed DNA-repair
pathway and propose a clinical confirmation.
\end{exercise}

\paragraph{Solution sketch.} Microsatellite instability is the
characteristic signature of mismatch-repair (MMR) failure: short
repeat tracts (di-, tri-, and tetra-nucleotide microsatellites)
contract or expand because the slippage errors introduced by DNA
polymerase at repeat sequences are no longer corrected. The
candidate failed pathway is MMR, with the most common molecular
mechanisms being MLH1 promoter methylation (the leading sporadic
cause, accounting for $\sim 75\%$ of microsatellite-unstable colorectal
cancers) or germline mutation in MLH1, MSH2, MSH6, or PMS2 (Lynch
syndrome). Clinical confirmation: (1) immunohistochemistry on tumour
tissue for the four MMR proteins, identifying which is absent;
(2) MSI PCR or MSI-by-NGS to confirm the microsatellite phenotype;
(3) MLH1 promoter-methylation testing to distinguish sporadic from
germline cases; (4) if methylation-negative, germline sequencing of
the MMR genes from blood for Lynch-syndrome diagnosis. The clinical
implication is double: pembrolizumab and other immune checkpoint
inhibitors have unusually high response rates in MSI-high tumours,
and Lynch-syndrome confirmation triggers cascade testing of
first-degree relatives.

\begin{exercise}
A bacterial strain in industrial fermentation drifts in its
production yield over many batches. Identify three possible
information-level causes (plasmid instability, mutation in the
production gene, regulatory drift, host evolution) and propose a
diagnostic measurement for each.
\end{exercise}

\paragraph{Solution sketch.} (1) Plasmid loss or plasmid copy-number
drift: diagnose by qPCR of a plasmid-specific sequence against a
chromosomal reference, comparing high-producing and low-producing
batches. Loss of plasmid in a fraction of cells reduces effective
production capacity. (2) Mutation in the production gene (cassette):
amplicon-sequence the production-gene region from the affected batch
and compare to the master seed bank. Loss-of-function or activity-
reducing mutations in the production gene or its promoter accumulate
under selection pressure if the production load is fitness-costly.
(3) Regulatory drift, including silencing of the inducible promoter:
RT-qPCR of the production-gene mRNA in the affected batch versus the
master seed bank; flow-cytometry of a promoter-GFP fusion to detect
silenced sub-populations. (4) Host evolution: whole-genome sequencing
of the affected batch versus the master seed bank, looking for
suppressor mutations in the host chromosome that reduce production-
related fitness costs. The systematic diagnostic plan distinguishes
the four causes; the engineering response (plasmid stabilisation by
antibiotic-free selection, chromosomal integration of the production
cassette, redesign of the regulatory architecture, or seed-bank
discipline) depends on which cause is identified.

\begin{exercise}
A child's whole-genome sequencing identifies $50$ new mutations
relative to the parents, of which two are in protein-coding regions.
Estimate the probability that one of the two coding-region mutations
is pathogenic and propose a follow-up workflow.
\end{exercise}

\paragraph{Solution sketch.} $50$ de novo mutations is the expected
range from Kong et al. \cite{paper:v6c03-kong-2012}; $2$ in coding
regions is in line with the $\sim 1.5\%$ coding fraction of the
genome. Per coding mutation, $\sim 70\%$ are missense, $\sim 25\%$
synonymous, $\sim 5\%$ nonsense or frameshift. The pathogenicity
prior is dominated by gene-level constraint: known haploinsufficient
genes (those whose loss of one functional copy produces disease,
e.g.\ \emph{SCN1A}, \emph{KMT2D}) raise the pathogenicity probability
of a deleterious mutation to $> 10\%$; mutations in low-constraint
genes have pathogenicity probability $< 0.1\%$. Workflow: (1) annotate
each mutation against gnomAD allele frequencies (a variant present
in unrelated healthy individuals is unlikely to be a strongly
penetrant de novo cause), ClinVar pathogenicity records, and
gene-level constraint metrics (pLI, missense Z); (2) for missense
variants, compute computational pathogenicity scores (CADD,
REVEL, AlphaMissense); (3) for variants in known disease genes,
correlate with the patient's clinical phenotype using a HPO-based
matching tool; (4) confirm any candidate by Sanger sequencing and
test the parents to confirm de novo status; (5) if no clear cause is
identified, return the case for trio re-analysis at six-month
intervals as gene-disease association improves. The expected
pathogenic yield from two coding de novo mutations in an unselected
population is in the $1$--$5\%$ range; in a developmentally affected
child it rises to $25$--$40\%$.

\subsection*{Failure analysis}\pathtag{core}

\begin{exercise}
A patient with xeroderma pigmentosum develops severe skin cancers
from minimal UV exposure. Identify the failed repair pathway
(nucleotide excision repair) and trace the chain from defective
repair to the observed cancer phenotype.
\end{exercise}

\paragraph{Solution sketch.} Xeroderma pigmentosum is caused by
biallelic loss-of-function mutations in one of seven nucleotide
excision repair (NER) genes (XPA through XPG) or in DNA polymerase
$\eta$ (XPV). NER is the pathway responsible for removing bulky DNA
adducts including UV-induced cyclobutane pyrimidine dimers (CPDs)
and $6$-$4$ photoproducts. With NER absent, every $\sim 10^4$ UV-
induced lesions per cell per hour of solar exposure persist into
the next round of replication; the replication fork encounters the
lesion and either stalls (triggering an alternative trans-lesion
synthesis pathway that is itself mutagenic) or generates a deletion
or substitution at the lesion position. The signature mutations
are C-to-T substitutions at pyrimidine dimers, the classic UV
mutational signature. With mutation rates elevated by approximately
three orders of magnitude in skin cells, the multi-hit cancer
progression (Section~3.7) accelerates dramatically: melanoma and
non-melanoma skin cancers appear in childhood rather than late
adulthood. The chain from defective repair to cancer phenotype is
NER loss $\to$ persistent CPDs $\to$ replicative mutation at each
CPD $\to$ accumulation in driver genes (CDKN2A, TP53, BRAF, KIT)
$\to$ proliferative dysregulation $\to$ tumour. The lifetime
incidence of skin cancer in XP patients exceeds $90\%$; clinical
management is rigorous sun avoidance from infancy.

\begin{exercise}
A new biopharmaceutical produced in mammalian cell culture shows
unexpected sequence variants at $0.5\,\%$ frequency in mass-spectrometric
analysis of the product. Identify three possible sources (replication
error, transcription error, translation error, post-translational
modification) and propose a measurement that would distinguish them.
\end{exercise}

\paragraph{Solution sketch.} (1) Replication error in the production
cell line: would manifest as a heritable sequence variant present in
a sub-clone of the production culture; distinguish by genomic
sequencing of single-cell clones from the production bioreactor,
which would localise the variant to a specific sub-population.
(2) Transcription error by the RNA polymerase: would manifest as
heterogeneity at the mRNA level not the DNA level; distinguish by
RNA-seq of the production-gene mRNA, which would show variant
allele frequencies of $\sim 10^{-5}$ at random positions (transcription
error rate), too low to explain $0.5\%$ unless a hotspot is involved.
(3) Translation misincorporation at a specific codon (e.g.\ a rare
codon paired with a saturated tRNA pool causing near-cognate
misreading): would manifest as protein-level variants without
corresponding mRNA-level variants; distinguish by comparing RNA-seq
to mass-spectrometric proteomics. (4) Post-translational modification
(deamidation, oxidation, glycosylation drift): would shift the mass
by a characteristic increment; distinguish by tandem MS fragmentation
patterns. The $0.5\%$ frequency is too high for transcription error
and too low for typical cell-line replication drift, so the prior
should weight (3) and (4) most heavily; the protocol is RNA-seq
(rules out 1 and 2) plus careful PTM mapping by LC-MS/MS, with codon-
context analysis of the variant positions to identify suspect
rare-codon hotspots.

\subsection*{Judgment}\pathtag{standard}

\begin{exercise}
A direct-to-consumer genetic-testing company reports a customer's
risk of a complex disease based on a polygenic-risk score. Identify
two engineering checks the customer should apply to the reported
probability (population calibration, ancestry-specific accuracy,
confidence interval) and write a one-paragraph reply to the customer's
question ``does this mean I will get the disease?''.
\end{exercise}

\paragraph{Discussion (rubric).} A good answer identifies three or
more of: (a) the reported number is a population-relative score, not
an absolute probability, and it changes meaning under cohort
reweighting; (b) most polygenic risk scores in 2024 are trained on
European-ancestry samples and have substantially lower predictive
accuracy in other ancestries (often $50$--$80\%$ less variance
explained); (c) the confidence interval on the score is wide (often
$\pm 30\%$ of the reported value or larger) because the per-variant
effect-size estimates carry estimation uncertainty; (d) the lifetime
prevalence of the disease determines the baseline against which the
score is read (a $2 \times$ increase on a $1\%$ baseline is far less
serious than a $2 \times$ increase on a $20\%$ baseline); (e) modifiable
risk factors (lifestyle, environment) often have comparable or larger
effects than the score. A high-credit reply might read:
``This score indicates that, relative to the population the company
trained the score on, you are estimated to be at a somewhat higher
genetic-component risk than average. It is not an individual
probability; it is a population-relative score with substantial
uncertainty. The score is most accurate for people of the ancestry
group that the training population reflected; for other ancestries
the predictive accuracy is reduced. The score does not account for
lifestyle factors or environment, which often contribute as much
or more risk than the genetic component. The score is one input to
a clinical-risk evaluation; consultation with a clinician who can
combine the score with your family history and modifiable risk
factors will give a more informative estimate.''

\begin{exercise}
A press release describes a ``DNA-based diagnostic'' that detects
disease ``before any symptoms appear'' with $99.9\,\%$ accuracy.
Identify the engineering questions (sensitivity, specificity, base
rate, positive predictive value) the accuracy claim does not answer
and write a one-paragraph reply.
\end{exercise}

\paragraph{Discussion (rubric).} A good answer separates the three
components of test performance: (a) sensitivity (P(positive $\mid$
disease)) and specificity (P(negative $\mid$ no disease)) are
independent of population prevalence and are what a properly designed
validation study reports; ``$99.9\%$ accuracy'' obscures the trade-off
between the two; (b) positive predictive value (P(disease $\mid$
positive)) and negative predictive value depend on the prevalence in
the screened population and are what the patient experiences; (c)
for a rare disease (prevalence $\sim 0.1\%$), a test with $99.9\%$
specificity gives PPV $\approx 50\%$, so half of all positives are
false; a test with $99\%$ specificity gives PPV $\approx 10\%$; (d)
the validation study's reference cohort determines the realism of
the claim, and a press release that does not name the validation
cohort, the prevalence in that cohort, and the separated sensitivity
and specificity has not produced a usable engineering claim. A
high-credit reply applies these to the specific test.

\begin{exercise}
A textbook describes the genome as a ``blueprint'' for the organism.
Identify two ways in which the blueprint metaphor is accurate
(sequence-defined inheritance, deterministic transcription rules) and
two ways in which it is misleading (the genome does not specify a
unique adult organism, and the regulatory layer is at least as
important as the coding sequence).
\end{exercise}

\paragraph{Discussion (rubric).} Accurate parts of the metaphor: the
genome is a sequence-defined, heritable, digital record of the protein
parts list, and the transcription-translation machinery operates by
deterministic rules that map sequence to product. Misleading parts:
(a) the genome does not encode a unique adult organism; the same
genome in different developmental, environmental, or stochastic
contexts produces different cell types and organism states; (b) the
regulatory layer (transcription factors, enhancers, epigenome,
chromatin) is dynamical and is at least as important as the coding
sequence for cell-state specification, and the regulatory layer is
not itself a fixed blueprint; (c) the developmental program is an
iterated, time-extended process that the genome influences but does
not specify in detail; (d) the blueprint metaphor invites the false
inference that an annotated genome ``determines'' the organism, when
in fact organism-level phenotype prediction from genome remains an
unsolved problem for any complex trait. A high-credit answer uses
the metaphor's failure to motivate the developmental and regulatory
content of the rest of the volume.

\begin{exercise}
A vendor advertises a ``synthetic genome'' that can be installed in
an empty cellular chassis to produce a fully functional cell.
Identify the engineering constraints the design must satisfy
(replication-origin compatibility, codon usage matched to host tRNA
pool, essential-gene coverage, regulatory completeness) and write a
one-paragraph reply identifying which constraint is most likely
underspecified in the advertisement.
\end{exercise}

\paragraph{Discussion (rubric).} The constraints any synthetic-genome
project must satisfy: (a) every essential gene must be present and
functional (the JCVI Mycoplasma syn3.0 minimal cell carries $473$
genes, of which approximately $149$ have no annotated function);
(b) the genome must carry replication origins compatible with the
chassis's replication machinery; (c) the codon usage must match the
host tRNA pool sufficiently to support translation; (d) the regulatory
architecture must be complete (every essential gene must be expressed
at the right time and level); (e) metabolic intermediate supplies
must match the chassis's biosynthesis capacity. The constraint most
often underspecified in commercial advertisements is (a) combined
with (d): claims of ``functional'' synthetic cells often rest on
genomes that carry the essential parts list but not the complete
regulatory architecture, so the cell grows but exhibits abnormal
growth dynamics, sensitivity to nutrient state, or genome instability.
A high-credit reply asks for the cell's doubling time, the genome's
stability over $100$ generations, the temperature- and nutrient-
robustness profile, and the comparison to the chassis's parent
strain on standard physiological assays.

\begin{exercise}
A press release claims that a new gene-editing therapy ``corrects''
a disease mutation in $99\,\%$ of patient cells. Identify the
engineering questions the percentage does not answer (which tissues
were edited, what fraction of stem cells were edited, what off-target
edits occurred, whether the correction is heritable) and write a
one-paragraph reply.
\end{exercise}

\paragraph{Discussion (rubric).} The $99\%$ figure is a single
aggregate and obscures the relevant engineering distinctions:
(a) \emph{which cells}: is the $99\%$ over the tissue of interest
(stem cells or the proliferating compartment), over circulating cells
(which turn over in months and so do not represent durable correction),
or over all cells in the body? (b) \emph{which sub-population}:
stem cells set the long-term outcome; $99\%$ correction of
mature cells with $1\%$ correction of stem cells produces transient
benefit, the failure mode of early gene therapies; (c) \emph{off-target
editing}: was the $1\%$ that was not corrected un-edited (safe), or
was it edited at the wrong locus (potentially deleterious)?
(d) \emph{heritability}: is the editing somatic only or has the
germline been affected? somatic editing is the current standard and
heritable editing is the He Jiankui case
\cite{acc:he-jiankui-crispr-2018}; the distinction is determinative
for the regulatory category; (e) \emph{durability}: how long after
treatment was the $99\%$ measured? a peak at $1$ month is not a
durable correction. A high-credit reply asks for the editing
breakdown by cell type and time, the off-target characterisation, and
the long-term follow-up plan.

\begin{exercise}
A reader argues that the discovery of the genetic code in the
$1960$s exhausted the engineering content of molecular biology.
Write a one-paragraph reply naming two layers of regulatory
information (chromatin state, non-coding RNA networks) that the
code does not capture and that have emerged as engineering targets
in the intervening decades.
\end{exercise}

\paragraph{Discussion (rubric).} The genetic code answers the question
``how does sequence specify protein?'' but leaves untouched several
layers of biological information that have become engineering targets
since. The reader should identify two or more of: (a) chromatin state
(DNA methylation, histone modifications, three-dimensional genome
organisation) carries cell-type-specific gene-expression information
that is heritable across cell divisions without sequence change;
(b) non-coding RNAs (microRNAs, long non-coding RNAs, circRNAs) form
regulatory networks that adjust mRNA stability, translation
efficiency, and gene-expression timing on scales of minutes to hours;
(c) alternative splicing produces multiple protein variants from a
single gene; the splice-site sequence rules are not contained in the
genetic code; (d) the gene-regulatory-network architecture
(transcription-factor circuits, signalling cascades) is the engineering
substrate of cell-state specification and developmental programming;
(e) the protein-folding code (the mapping from sequence to
three-dimensional structure) is itself unsolved by the genetic code
and was a major open problem until AlphaFold (2021). Each of these
layers has emerged as a distinct engineering and therapeutic
target since the 1960s, justifying the claim that molecular biology
remains an active engineering discipline.
